% !TeX program = xelatex
\documentclass[11pt,openany]{book}
\usepackage{fontspec}
\setmainfont{Latin Modern Roman}
\setsansfont{Latin Modern Sans}
\setmonofont{Latin Modern Mono}
\usepackage[english]{babel}
\usepackage[a4paper,margin=2.45cm]{geometry}
\usepackage{amsmath,amssymb,amsfonts,bm,mathtools}
\usepackage{booktabs,array,longtable,tabularx}
\usepackage{amsthm}
\usepackage{tikz}
\usepackage{pgfplots}
\pgfplotsset{compat=1.18}
\usepackage{fancyhdr}
\usepackage{placeins}
\usepackage{makeidx}
\makeindex
\usepackage{enumitem}
\usepackage{hyperref}
\usepackage{xcolor}
\usepackage[most]{tcolorbox}
\usepackage{microtype}
\usepackage{needspace}

\hypersetup{
  colorlinks=true,
  linkcolor=blue,
  urlcolor=blue,
  citecolor=blue,
  pdftitle={A Concise Course of Computational Physics: From Dynamical Reduction to Modern High-Order Discretization},
  pdfauthor={}
}

\setcounter{tocdepth}{1} % Show parts, chapters, and sections.
\setlength{\parindent}{2em}
\setlength{\parskip}{0.35em}

\newcommand{\dd}{\,\mathrm d}
\newcommand{\ee}{\mathrm e}
\newcommand{\ii}{\mathrm i}
\newcommand{\R}{\mathbb R}
\newcommand{\Order}{\mathcal O}
\newcommand{\Lop}{\mathcal L}
\newcommand{\Nop}{\mathcal N}
\newcommand{\Ksp}{\mathcal K}
\newcommand{\avg}[1]{\left\langle #1\right\rangle}

\newtcolorbox{chaptersummary}{
  colback=blue!3,
  colframe=blue!45!black,
  title={Chapter Overview},
  fonttitle=\bfseries,
  boxrule=0.6pt,
  arc=2mm,
  left=2mm,
  right=2mm,
  top=1mm,
  bottom=1mm
}

\newtcolorbox{authornote}{
  colback=yellow!8,
  colframe=orange!75!black,
  title={Potential Conceptual Jump},
  fonttitle=\bfseries,
  boxrule=0.8pt,
  arc=2mm,
  left=2mm,
  right=2mm,
  top=1mm,
  bottom=1mm
}

\newtcolorbox{askai}{
  colback=violet!3,
  colframe=violet!55!black,
  title={Ask AI},
  fonttitle=\bfseries,
  boxrule=0.7pt,
  arc=2mm,
  left=2mm,
  right=2mm,
  top=1mm,
  bottom=1mm
}

\newcommand{\studyroute}[3]{%
  \par\smallskip\noindent\textbf{Further Detail and Search. }#1\par
  \noindent\textbf{Keywords: }#2.\par
  \noindent\textbf{Ask AI: }#3\par
}


\emergencystretch=2em
\raggedbottom
\setlist{nosep,topsep=3pt}
\pagestyle{fancy}
\fancyhf{}
\fancyhead[LE]{\small A Concise Course of Computational Physics: Volume I}
\fancyhead[RO]{\small\nouppercase{\rightmark}}
\fancyfoot[C]{\thepage}
\setlength{\headheight}{15pt}
\newcommand{\learninggoals}[1]{\paragraph{Learning Objectives}#1}
\newcommand{\conceptbridge}[3]{%
\begin{authornote}
\textbf{Potential Question: }#1\par
\textbf{Author Supplement: }#2\par
\textbf{Core Conceptual Connection: }#3
\end{authornote}}
\newcommand{\aiexercise}[2]{%
\begin{askai}
\textbf{Ask AI: }#1\par\smallskip
\textbf{Follow-up: }#2
\end{askai}}
\newcommand{\norm}[1]{\left\lVert#1\right\rVert}
\newcommand{\jump}[1]{\left[#1\right]}
\newcommand{\avgface}[1]{\left\{#1\right\}}
\newcommand{\Poly}{\mathbb P}
\newcommand{\tr}{\operatorname{tr}}
\newcommand{\diag}{\operatorname{diag}}
\numberwithin{equation}{chapter}
\newtheorem{theorem}{Theorem}[chapter]
\newtheorem{proposition}[theorem]{Proposition}
\theoremstyle{definition}
\newtheorem{definition}[theorem]{Definition}

\title{From Dynamical Reduction\\to Modern High-Order Discretization}
\author{Harvey Tan\\[0.6em]\large A Concise Course of Computational Physics\\Volume I}
\date{November 2026}
\hypersetup{pdfauthor={Harvey Tan}}
\begin{document}
\frontmatter
\maketitle
\chapter*{Preface}
\addcontentsline{toc}{chapter}{Preface}
\begin{quote}
\itshape
This book is dedicated to Landau and Lifshitz.\\
With gratitude for guiding generations to the beauty of theory and the wonder it inspires.
\end{quote}
\begin{tcolorbox}[colback=red!3,colframe=red!65!black,title={Required: AI Explanations and Code Experiments},fonttitle=\bfseries,boxrule=0.8pt,arc=1mm,breakable]
\textbf{Do not study from the text alone. Use AI explanations and code experiments alongside these books.}
The text is a highly compressed knowledge map. When a conceptual link, derivation,
or prerequisite is unclear, give AI the current notation, assumptions, and precise
question and request a step-by-step explanation. Run the companion code or a minimal
example, vary parameters, and inspect error, conservation, or stability.
Check AI explanations, derivations, and code against the stated assumptions,
references, and numerical results.
\end{tcolorbox}
Computational physics represents physical states, reduces degrees of freedom, closes unresolved effects, and constructs finite-dimensional evolution under finite computational resources. Each step must identify the information retained, the assumptions introduced, and the structures that can be verified. This volume evaluates these choices through conservation, stability, symmetry, and wave errors.

Four questions organize the text: how to represent the state; how scales and statistics reduce dynamics; how physical assumptions close the influence of unresolved variables; and how finite-dimensional residuals are constrained by discretization. Model selection determines the computational object; numerical methods determine its approximate representation and evolution. Structure preservation connects these decisions.

Part~I proceeds from Hamiltonian dynamics to Liouville, kinetic, and continuum descriptions. Part~II builds function spaces and weak formulations through integration by parts, comparing trial spaces, test spaces, and residual constraints. Part~III derives DG, FR, VCJH, mortar coupling, and time integration. Parts~IV and~V develop discrete-velocity and physical-reduction routes. Part~VI connects meshes, method selection, and verification.

The intended reader has upper-undergraduate mathematical preparation: multivariable calculus, linear algebra, ODEs, PDEs, and basic numerical analysis. The main text supplies key definitions, boundary terms, and algorithmic constructions. General trace theorems, error theory, and research topics have entry points through further-detail notes, search terms, and specific AI questions. Learning objectives and conceptual connections locate each chapter within the dependency chain.

At the formation of algebraic systems, this volume points to Volume~0, \emph{Numerical Linear Algebra and Modern High-Performance Computing}, which develops matrix structure, factorization, Krylov methods, preconditioning, multigrid, and parallel execution. Volume~II uses Maxwell's equations to study geometry, constraint propagation, and electromagnetic discretization.

\paragraph{Notation}
$\rho_N$ is a normalized many-particle probability density, $f_s$ a number-density distribution, $\rho$ a mass density, and $\rho_c$ a charge density. The scalar field is $u$; $U$ denotes conserved variables or their discrete coefficients, with dimensions specified for each model. Mesh size, polynomial degree, and time step are $h$, $p$, and $\Delta t$. In Chapters~10--14, one-dimensional interfaces use left trace $u^-$, right trace $u^+$, jump $[u]=u^+-u^-$, and average $\{u\}=(u^-+u^+)/2$. Normal numerical fluxes additionally include normal orientation. Collision times, viscous stresses, and discrete relaxation times are distinguished within each chapter.

\tableofcontents
\mainmatter
\part{From Microscopic Dynamics to Continuum Models}
\chapter{From Many-Particle Dynamics to Kinetic Equations}\label{ch:1}

\begin{chaptersummary}
Moving from particle trajectories to computable models requires choosing which information to retain. Probability density represents uncertainty in the microscopic state, marginal integration reduces explicit coordinates, and closure approximates higher correlations entering lower-order evolution. This reduction chain provides the starting point for continuum and plasma models.
\end{chaptersummary}
\learninggoals{Relate trajectory and probability descriptions; distinguish marginal integration, correlation hierarchies, and collision closure; check distribution normalizations and conserved moments.}
\studyroute{The term-by-term integration from the Liouville equation to the BBGKY hierarchy and the limiting construction of collision operators from two-particle correlations are omitted.}{Liouville equation; marginal distribution; BBGKY hierarchy; molecular chaos}{For two particles interacting through pair forces, define the one-particle marginal by integration and derive the term in its equation that depends on the two-particle distribution.}

\section{Hamiltonian Dynamics and Phase-Space Transport}\index{Liouville equation}
Consider \(N\) particles, with the position and momentum of particle \(i\) denoted by
\[
\bm x_i,\qquad \bm p_i.
\]
The full system state is a point in \(6N\)-dimensional phase space:
\[
\Gamma=
(\bm x_1,\dots,\bm x_N,\bm p_1,\dots,\bm p_N).
\]

For a Hamiltonian \(H\), the equations of motion are
\[
\dot{\bm x}_i=\frac{\partial H}{\partial \bm p_i},
\qquad
\dot{\bm p}_i=-\frac{\partial H}{\partial \bm x_i}.
\]

When macroscopic statistics are the quantities of interest and particle-by-particle computation is too expensive, introduce a phase-space probability density
\[
\rho_N(\Gamma,t).
\]

\begin{authornote}
\textbf{Potential Question:}
"Why move from deterministic \(N\)-particle trajectories to a probability density?"

\textbf{Author Supplement:}
Many engineering problems concern density, temperature, or mean flow rather than individual trajectories. For systems containing enormous numbers of particles, tracking every degree of freedom is usually prohibitively expensive; direct evaluation of all pair interactions alone costs $O(N^2)$. A probability description focuses on possible states and their statistical laws, and represents uncertainty in the initial microscopic state. Particle methods or Eulerian--Lagrangian hybrid methods remain available when particular particles must be tracked.

\textbf{Core Conceptual Connection:}
\[
\text{deterministic particle trajectory}
\longrightarrow
\text{phase-space probability description}
\]
\end{authornote}

Conservation of probability gives a high-dimensional continuity equation:
\[
\frac{\partial \rho_N}{\partial t}
+
\nabla_\Gamma\cdot(\rho_N\dot\Gamma)=0.
\]

\begin{askai}
Ask AI:
\begin{quote}
Starting from the ordinary three-dimensional continuity equation, extend the analogy to \(6N\)-dimensional phase space and explain step by step why \(\rho_N(\Gamma,t)\) satisfies
\[
\partial_t\rho_N+\nabla_\Gamma\cdot(\rho_N\dot\Gamma)=0.
\]
Explain explicitly what flux means in phase space.
\end{quote}

Follow-up:
\begin{quote}
Use a Hamiltonian system with one position and one momentum coordinate to show how the phase-space velocity field transports a probability cloud.
\end{quote}
\end{askai}

Hamiltonian systems satisfy
\[
\nabla_\Gamma\cdot\dot\Gamma=0,
\]
Hence
\[
\frac{\partial \rho_N}{\partial t}
+
\dot\Gamma\cdot\nabla_\Gamma\rho_N=0,
\]
or
\[
\frac{\partial \rho_N}{\partial t}+\{\rho_N,H\}=0.
\]

The Liouville equation describes transport of probability density by microscopic dynamics; Hamiltonian flow preserves phase-space volume. The equation remains \(6N\)-dimensional. Only subsequent marginalization reduces the number of explicit coordinates.

\begin{authornote}
\textbf{Potential Question:}
"Why does a probability description remain high-dimensional, and which information does marginalization retain or discard?"

\textbf{Author Supplement:}
The density \(\rho_N\) depends on all \(6N\) phase-space coordinates, so a probability description alone does not reduce dimension. Even the one-particle distribution \(f(\bm x,\bm v,t)\) requires approximately \(N_x^3N_v^3\) degrees of freedom on a direct tensor-product grid: every spatial point carries a three-dimensional velocity distribution. With 100 points per coordinate direction, there are \(10^{12}\) values; one double-precision scalar array alone occupies approximately 8 TB. Actual cost also depends on dimensional reduction, mesh design, collision model, and parallel implementation.

Marginalization integrates out the other particles' explicit positions, velocities, and joint information unavailable in a one-particle distribution. Interactions still transmit the effects of these correlations to the retained variables, motivating closure. The roles are distinct: marginalization defines retained variables, closure models unresolved correlations, and initial and boundary conditions supply space-time data.

\textbf{Core Conceptual Connection:}
\[
\rho_N(\Gamma,t)
\longrightarrow
f(\bm x,\bm v,t)
\]
\end{authornote}

\section{Marginal Distributions, Normalization, and Correlation Hierarchies}\index{Marginalization}
For particles of the same mass $m$, replace momentum coordinates by velocity coordinates and define
\[
\varrho_N(\bm x_1,\bm v_1,\ldots,\bm x_N,\bm v_N,t)
=m^{3N}\rho_N(\bm x_1,\ldots,m\bm v_N,t),\qquad
\int\varrho_N\dd z_1\cdots\dd z_N=1,
\]
Here $z_i=(\bm x_i,\bm v_i)$, and every momentum argument of $\rho_N$ on the right is evaluated at $\bm p_i=m\bm v_i$.
Probability-normalized and number-density marginals are, respectively,
\[
g_s(z_1,\ldots,z_s)=\int\varrho_N\dd z_{s+1}\cdots\dd z_N,\qquad
f_s=\frac{N!}{(N-s)!}g_s.
\]
Thus $\int g_s=1$, $\int f_1=N$, and $\int f_2=N(N-1)$. Below, $f=f_1$.
Marginalization is performed at each fixed time; the evolution of the retained variables follows by integrating the Liouville equation.

Define the one-particle distribution
\[
f(\bm x,\bm v,t),
\]
so that
\[
f(\bm x,\bm v,t)\dd\bm x\dd\bm v
\]
is the expected number of particles in the corresponding position and velocity elements.

Integrating the \(N\)-particle equation over the other variables produces
\[
f_1\longrightarrow f_2\longrightarrow f_3\longrightarrow\cdots,
\]
the BBGKY hierarchy. The one-particle equation depends on two-particle correlations, and the two-particle equation depends on three-particle correlations. Closing the model requires assumptions about these correlations.

\begin{authornote}
\textbf{Potential Question:}
"Why does marginalization produce a hierarchy \(f_1,f_2,f_3,\dots\), rather than a closed one-particle equation?"

\textbf{Author Supplement:}
The mean force on a selected particle depends on the probability that it and another particle occupy particular states simultaneously: the two-particle distribution. Integration removes explicit coordinates but not interactions. Similarly, the evolution of the two-particle distribution involves the three-particle distribution. Solving only lower-order equations therefore requires approximating higher-order correlations under stated conditions. This step is called closure.

\textbf{Core Conceptual Connection:}
\[
\text{marginalization}
\longrightarrow
\text{BBGKY hierarchy}
\longrightarrow
\text{closure problem}
\]
\end{authornote}

\section{Kinetic Equations and Collision Models}\index{Boltzmann equation}
A general kinetic equation has the form
\[
\frac{\partial f}{\partial t}
+\bm v\cdot\nabla_{\bm x}f
+\frac{\bm F}{m}\cdot\nabla_{\bm v}f
=C[f].
\]
The left side describes phase-space transport; the right side describes collisions or relaxation. The displayed form assumes $\nabla_{\bm v}\cdot\bm F=0$, as for velocity-independent forces and the Lorentz force. A general velocity-dependent force requires the velocity-space divergence form.

Neglecting discrete collisions gives the Vlasov equation for charged particles:
\[
\frac{\partial f_s}{\partial t}
+\bm v\cdot\nabla_{\bm x}f_s
+\frac{q_s}{m_s}
(\bm E+\bm v\times\bm B)\cdot\nabla_{\bm v}f_s=0.
\]
Charge and current densities are given by
\[
\rho_c=\sum_s q_s\int f_s\dd\bm v,
\qquad
\bm J=\sum_s q_s\int \bm v f_s\dd\bm v
\]
and feed back into Maxwell's equations, forming the self-consistent Vlasov--Maxwell system.

For a dilute neutral gas, approximating collisions as instantaneous binary encounters gives the Boltzmann equation:
\[
\frac{\partial f}{\partial t}
+\bm v\cdot\nabla_{\bm x}f
+\frac{\bm F}{m}\cdot\nabla_{\bm v}f
=Q(f,f),
\]
where
\[
Q(f,f)
=
\int B\left(f'f_*'-ff_*\right)\dd\Omega\dd\bm v_*.
\]
Apply molecular chaos to incoming particles, using the normalization of a local collision pair:
\[
f_2(\bm v,\bm v_*)\approx f(\bm v)f(\bm v_*).
\]

For the finite-$N$ number-density normalization above, the exact product for independent particles also contains a factor $(N-1)/N$; the usual collision closure absorbs this factor in the large-particle-number limit. The assumption of uncorrelated particles applies to incoming-particle statistics in the corresponding dilute limit.

\begin{askai}
Ask AI:
\begin{quote}
Use a minimal example with two particle velocities to explain \(f_2(\bm v,\bm v_*)\approx f(\bm v)f(\bm v_*)\). Distinguish the meanings of "uncorrelated" in probability theory, collision physics, and closure of the Boltzmann equation.
\end{quote}

Follow-up:
\begin{quote}
Construct a counterexample with \(f_2\neq f f_*\), and identify the information absent from the one-particle Boltzmann equation.
\end{quote}
\end{askai}

The full collision integral is expensive. The BGK model replaces it with a single relaxation process:
\[
C[f]=\frac{f^{\mathrm{eq}}-f}{\tau},
\]
where
\[
f^{\mathrm{eq}}
=
n\left(\frac{m}{2\pi k_BT}\right)^{3/2}
\exp\left[
-\frac{m|\bm v-\bm u|^2}{2k_BT}
\right].
\]

For long-range Coulomb interactions in plasmas, particles typically undergo many small-angle deflections. This leads to Landau or Fokker--Planck operators, which describe collisions as velocity-space drift and diffusion.

Different collision assumptions therefore define a model hierarchy:
\[
\begin{gathered}
\text{Collisionless Vlasov}
\longrightarrow\\
\text{Weakly collisional Vlasov--Fokker--Planck}\\
\longrightarrow
\text{Strongly collisional fluid model}.
\end{gathered}
\]

\chapter{From Velocity Moments to Continuum Equations}\label{ch:2}

\begin{chaptersummary}
The distribution function in the preceding chapter retains the full velocity dependence. This chapter defines density, mean velocity, pressure, and energy through velocity moments, derives macroscopic conservation equations, and supplies unresolved stress and heat flux through constitutive relations. This provides the basis for Euler and Navier--Stokes models and finite-volume conservation. The retained moments and closure assumptions jointly determine the range of validity.
\end{chaptersummary}
\learninggoals{Derive conserved moments and identify unclosed terms; distinguish the pressure tensor, viscous stress, and heat flux; state the assumptions behind Euler and Navier--Stokes models.}
\studyroute{The full algebra of the energy moment and kinetic calculations of transport coefficients are omitted.}{velocity moments; Chapman--Enskog expansion; constitutive closure}{Starting from the three classes of conserved moments of $f$, distinguish exact moment identities from the additional assumptions in Newtonian viscosity and Fourier heat flux.}

\section{Conserved Moments and Macroscopic Observables}\index{Moment closure}
Consider a three-dimensional, single-species, monatomic dilute gas. Elastic collisions conserve mass, momentum, and energy; the external force is $m\bm b(\bm x,t)$. Assume sufficient decay of $f$ for velocity-space integration.
A distribution contains the full velocity dependence; a macroscopic model selects finitely many weighted integrals.

Number density, mass density, and macroscopic velocity are
\[
n=\int f\dd\bm v,
\qquad
\rho=mn,
\qquad
\rho\bm u=m\int \bm v f\dd\bm v.
\]

Define the peculiar velocity
\[
\bm c=\bm v-\bm u.
\]
The pressure tensor and heat flux are
\[
\bm P=m\int \bm c\bm c f\dd\bm v,
\qquad
\bm q=\frac{m}{2}\int |\bm c|^2\bm c f\dd\bm v.
\]

Taking the zeroth velocity moment and using particle-number conservation in collisions gives
\[
\int Q(f,f)\dd\bm v=0,
\]

\begin{authornote}
\textbf{Potential Question:}
"Why do velocity moments produce macroscopic equations? What do the zeroth and first moments measure physically?"

\textbf{Author Supplement:}
Density, mean velocity, and energy are weighted integrals of the full velocity distribution. Their direct connections to conservation laws, measurements, and macroscopic states make them natural continuum variables. A finite moment set also sets the descriptive resolution: different distributions with identical retained moments represent the same macroscopic state.

\textbf{Core Conceptual Connection:}
\[
f(\bm x,\bm v,t)
\longrightarrow
\text{moments}
\longrightarrow
\text{macroscopic observables}
\]
\end{authornote}

This gives the continuity equation:
\[
\boxed{
\frac{\partial\rho}{\partial t}
+\nabla\cdot(\rho\bm u)=0
}.
\]

Its control-volume form is
\[
\frac{\dd}{\dd t}\int_\Omega \rho\dd V
=
-\int_{\partial\Omega}\rho\bm u\cdot\bm n\dd S.
\]
which is the starting point of finite-volume methods.

The first moment gives the momentum equation:
\[
\frac{\partial(\rho\bm u)}{\partial t}
+
\nabla\cdot(\rho\bm u\bm u+\bm P)
=
\rho\bm b.
\]

\begin{askai}
Ask AI:
\begin{quote}
Starting from the general kinetic equation, take the first velocity moment step by step and derive
\[
\partial_t(\rho\bm u)+\nabla\cdot(\rho\bm u\bm u+\bm P)=\rho\bm b.
\]
Explain every integration by parts, the required decay at infinity, and the physical meaning of each term.
\end{quote}

Follow-up:
\begin{quote}
Decompose \(\bm P=m\int \bm c\bm c f\dd\bm v\) into isotropic pressure and deviatoric stress, and explain why it is not simply a scalar.
\end{quote}
\end{askai}

Writing
\[
\bm P=p\bm I-\bm\tau,
\]
gives
\[
\frac{\partial(\rho\bm u)}{\partial t}
+
\nabla\cdot(\rho\bm u\bm u)
=
-\nabla p+\nabla\cdot\bm\tau+\rho\bm b.
\]

\section{Constitutive Closure and Continuum Models}
Closure requires a constitutive relation between stress and macroscopic variables. A Newtonian fluid satisfies
\[
\bm\tau
=
\mu\left[\nabla\bm u+(\nabla\bm u)^T\right]
+\lambda(\nabla\cdot\bm u)\bm I.
\]

\begin{authornote}
\textbf{Potential Question:}
"Why is the momentum equation not the end of the derivation? Why does taking moments not automatically supply a constitutive law?"

\textbf{Author Supplement:}
Lower-order moment equations generally contain higher-order moments, including stress in the momentum equation and heat flux in the energy equation. A model with finitely many macroscopic unknowns must approximate these unresolved quantities using the retained variables and their gradients. Constitutive relations supply this additional physical assumption; their form and validity depend on the medium and the scales under study.

\textbf{Core Conceptual Connection:}
\[
\text{moment equations}
\longrightarrow
\text{unclosed stresses and heat fluxes}
\longrightarrow
\text{constitutive closure}
\]
\end{authornote}


For incompressible flow with constant density and viscosity,
\[
\nabla\cdot\bm u=0,
\]
\[
\rho\left(
\frac{\partial\bm u}{\partial t}
+\bm u\cdot\nabla\bm u
\right)
=
-\nabla p+\mu\nabla^2\bm u+\rho\bm b.
\]

The inertial term transports momentum, pressure redistributes it, and viscosity smooths velocity gradients. Neglecting viscosity and heat conduction gives the Euler equations.

The energy moment gives the total-energy equation:
\[
\frac{\partial(\rho E)}{\partial t}
+
\nabla\cdot[(\rho E+p)\bm u]
=
\nabla\cdot(\bm\tau\cdot\bm u)
-\nabla\cdot\bm q
+\rho\bm b\cdot\bm u,
\]
Here $e$ is internal energy per unit mass. A monatomic ideal gas satisfies $p=\rho k_BT/m$ and $e=3k_BT/(2m)$. Newtonian stress and Fourier heat flux correspond to near-local-equilibrium conditions and scales long compared with the mean free path.
\[
E=e+\frac12|\bm u|^2,
\qquad
\bm q=-k\nabla T.
\]

The transition from kinetics to a continuum model retains finitely many moments and closes higher-order moments through constitutive relations.

\chapter{Electromagnetic Fields, Maxwell's Equations, and Self-Consistent Coupling}\label{ch:3}

\begin{chaptersummary}
The preceding chapters constructed matter models from particles and distributions. Electromagnetic fields now couple self-consistently to that matter: Maxwell constraints express charge conservation, field--particle feedback produces nonlinearity, and linearization extracts the operator governing small disturbances. This local response underpins dispersion, stability, and the electromagnetic discretizations of Volume II.
\end{chaptersummary}
\learninggoals{Connect charge continuity to Maxwell constraints; identify the data exchanged in Vlasov--Maxwell coupling; construct a Fr\'echet derivative at a specified base state and state the conditions for plane-wave analysis.}
\studyroute{Specific material models and kinetic responses are omitted in favor of the common structure of linearized operators.}{Fr\'echet derivative; linear response; Vlasov linearization}{Find the tangent tensor for $\bm D=\varepsilon\bm E+\alpha|\bm E|^2\bm E$ and explain how the base-state direction enters the linearized equations.}

\section{Electromagnetic Conservation and Field--Particle Coupling}\index{Vlasov--Maxwell system}
Maxwell's equations are
\[
\nabla\cdot\bm E=\frac{\rho_c}{\varepsilon_0},
\qquad
\nabla\cdot\bm B=0,
\]
\[
\frac{\partial\bm B}{\partial t}+\nabla\times\bm E=0,
\]
\[
\frac{\partial\bm E}{\partial t}
-c^2\nabla\times\bm B
=
-\frac{\bm J}{\varepsilon_0}.
\]

Taking the divergence of the Amp\`ere--Maxwell equation and using Gauss's law gives
\[
\frac{\partial\rho_c}{\partial t}
+\nabla\cdot\bm J=0.
\]

\begin{askai}
Ask AI:
\begin{quote}
Starting from the Amp\`ere--Maxwell equation
\[
\partial_t\bm E-c^2\nabla\times\bm B=-\bm J/\varepsilon_0
\]
take the divergence and use Gauss's law to derive charge conservation. Explain the structural role of \(\nabla\cdot(\nabla\times\bm B)=0\).
\end{quote}

Follow-up:
\begin{quote}
Construct discrete operators that do not satisfy \(D_hC_h=0\), and show how they can generate nonphysical charge.
\end{quote}
\end{askai}

Charge conservation is therefore built into the continuous structure of Maxwell's equations. A discretization that breaks this structure can generate nonphysical charge.

The Lorentz force density in a charged continuum is
\[
\bm f_L=\rho_c\bm E+\bm J\times\bm B.
\]
Combining fluid equations with Ohm's law gives two-fluid models or single-fluid MHD.
For kinetic species, $f_s$ determines $\rho_c$ and $\bm J$, while the fields change $f_s$ through $q_s(\bm E+\bm v\times\bm B)$. This two-way dependence is self-consistent field--particle coupling. Single-fluid MHD additionally uses low-frequency, quasineutral, and scale-ordering assumptions to reduce particle and species-fluid descriptions to a common velocity, pressure, and magnetic field. Chapter~\ref{ch:18} gives the closure.

\section{Nonlinear Operators and Local Linearization}
Vacuum Maxwell equations are linear. Coupling them to media, carriers, fluids, or nonlinear constitutive laws generally produces a nonlinear system. Write the coupled model abstractly as
\[
\Nop(\bm U)=0.
\]

\begin{authornote}
\textbf{Potential Question:}
"How does writing the coupled Maxwell system as \(\Nop(\bm U)=0\) help analyze perturbations?"

\textbf{Author Supplement:}
Time derivatives, spatial derivatives, and nonlinear couplings can be collected into an operator acting on the unknown fields. Without changing the problem, this notation provides a common description of perturbations near a solution. The operator derivative describes the propagation of small changes and supplies the linear equation approximately solved at each Newton step.

\textbf{Core Conceptual Connection:}
\[
\text{coupled field equations}
\longrightarrow
\Nop(\bm U)=0
\longrightarrow
\text{linearized operator}
\]
\end{authornote}

Near a base state \(\bm U_0\), write
\[
\bm U=\bm U_0+\epsilon\bm u,
\qquad 0<\epsilon\ll1.
\]
For a sufficiently smooth operator,
\[
\Nop(\bm U_0+\epsilon\bm u)
=
\Nop(\bm U_0)
+\epsilon D\Nop[\bm U_0](\bm u)
+\frac{\epsilon^2}{2}
D^2\Nop[\bm U_0](\bm u,\bm u)
+\cdots.
\]
If the base state satisfies
\[
\Nop(\bm U_0)=0,
\]
retaining first-order terms gives
\[
\Lop\bm u=0,
\qquad
\Lop:=D\Nop[\bm U_0].
\]

Here \(D\Nop[\bm U_0]\) is the Fr\'echet derivative, also called the linearized, Jacobian, or tangent operator. It is the local linear approximation of the nonlinear operator at the base state.

For example, let a nonlinear medium satisfy
\[
\bm D=\bm D(\bm E),
\qquad
\bm B=\bm B(\bm H).
\]
Set
\[
\bm E=\bm E_0+\epsilon\bm e,
\qquad
\bm H=\bm H_0+\epsilon\bm h.
\]
The first-order constitutive expansion is
\[
\bm D(\bm E_0+\epsilon\bm e)
=
\bm D(\bm E_0)
+
\epsilon
\left.
\frac{\partial\bm D}{\partial\bm E}
\right|_{\bm E_0}\bm e
+\Order(\epsilon^2),
\]
\[
\bm B(\bm H_0+\epsilon\bm h)
=
\bm B(\bm H_0)
+
\epsilon
\left.
\frac{\partial\bm B}{\partial\bm H}
\right|_{\bm H_0}\bm h
+\Order(\epsilon^2).
\]

\begin{askai}
Ask AI:
\begin{quote}
Compare finite-dimensional Taylor expansion with the Fr\'echet derivative, and explain step by step why the nonlinear material relation
\[
\bm D=\bm D(\bm E),\qquad \bm B=\bm B(\bm H)
\]
produces a tangent material tensor near a base state. Use an anisotropic matrix example.
\end{quote}

Follow-up:
\begin{quote}
If \(\bm D(\bm E)\) is nonlinear, why does the perturbation equation contain \(\partial\bm D/\partial\bm E|_{\bm E_0}\) rather than the original \(\bm D\)?
\end{quote}
\end{askai}

Define the tangent material tensors
\[
\bm\varepsilon_{\mathrm{tan}}
=
\left.\frac{\partial\bm D}{\partial\bm E}\right|_{\bm E_0},
\qquad
\bm\mu_{\mathrm{tan}}
=
\left.\frac{\partial\bm B}{\partial\bm H}\right|_{\bm H_0}.
\]
For instantaneous, nondispersive constitutive laws, further assume a time-independent base state, so the tangent tensors are constant in time, and set the perturbation current to zero. The first-order perturbations then satisfy a linear Maxwell-type system:
\[
\nabla\times\bm e
=
-\bm\mu_{\mathrm{tan}}\frac{\partial\bm h}{\partial t},
\]
\[
\nabla\times\bm h
=
\bm\varepsilon_{\mathrm{tan}}\frac{\partial\bm e}{\partial t}.
\]

\section{Vlasov Perturbations and Wave Dispersion}
In the Vlasov--Maxwell system, choose an equilibrium $f_{s0}$, $\bm E_0=0$, and specified $\bm B_0$ satisfying the zeroth-order equations, and write
\[
f_s=f_{s0}+\epsilon f_{s1},
\qquad
\bm E=\epsilon\bm E_1,
\qquad
\bm B=\bm B_0+\epsilon\bm B_1.
\]
Neglecting \(\Order(\epsilon^2)\) terms gives
\[
\frac{\partial f_{s1}}{\partial t}
+\bm v\cdot\nabla_{\bm x}f_{s1}
+\frac{q_s}{m_s}
(\bm v\times\bm B_0)\cdot\nabla_{\bm v}f_{s1}
=
-\frac{q_s}{m_s}
(\bm E_1+\bm v\times\bm B_1)
\cdot\nabla_{\bm v}f_{s0}.
\]
The perturbation current is
\[
\bm J_1=\sum_s q_s\int \bm v f_{s1}\dd\bm v.
\]

For a closed constant-coefficient linear system with finitely many amplitude unknowns, or a kinetic system whose velocity response has already been integrated, use plane waves:
\[
\bm u_1(\bm x,t)
=
\widehat{\bm u}
\exp[\ii(\bm k\cdot\bm x-\omega t)],
\]

\begin{authornote}
\textbf{Potential Question:}
"Why use plane waves to analyze perturbations? Does this assume that every solution is a wave?"

\textbf{Author Supplement:}
A constant-coefficient operator linearized about a uniform background preserves each Fourier wave number, allowing separate modal analysis. Slowly varying coefficients can be frozen locally to probe local response with plane waves. Boundaries and global nonuniformity enter through the corresponding global problem.

\textbf{Core Conceptual Connection:}
\[
\text{linear perturbation equation}
\longrightarrow
\text{plane-wave ansatz}
\longrightarrow
\text{dispersion relation}
\]
\end{authornote}

Then
\[
\partial_t\mapsto-\ii\omega,
\qquad
\nabla\mapsto\ii\bm k.
\]
and the linear PDE system becomes
\[
\bm M(\omega,\bm k)\widehat{\bm u}=0.
\]

\begin{askai}
Ask AI:
\begin{quote}
Starting with the plane-wave substitution
\[
\partial_t\mapsto-\ii\omega,\qquad \nabla\mapsto\ii\bm k
\]
choose a simple linear PDE system and show step by step how it becomes
\[
\bm M(\omega,\bm k)\widehat{\bm u}=0.
\]
Explain why a nonzero amplitude requires \(\det\bm M(\omega,\bm k)=0\).
\end{quote}

Follow-up:
\begin{quote}
Derive the complete dispersion relation for one-dimensional acoustic waves or vacuum Maxwell plane waves.
\end{quote}
\end{askai}

The condition for a nonzero solution,
\[
\det\bm M(\omega,\bm k)=0
\]
is the dispersion relation.

For kinetic equations that retain continuous velocity variables, the first result is a velocity-space response problem. An analytic continuation, integration contour, and field equations must be specified before the corresponding dielectric response and dispersion condition can be constructed.

Perturbation theory constructs solutions order by order in a small parameter; stability analysis examines linear perturbation growth rates. The Fr\'echet derivative also generates the Jacobian used in nonlinear algebraic iteration, whose solution algorithms are developed in Volume~0.

\chapter{Scale Analysis and Asymptotic Reduction}\label{ch:4}

\begin{chaptersummary}
Kinetic, continuum, and electromagnetic models contain several scales and mechanisms. Selecting among them begins with nondimensionalization, which expresses relative magnitudes as scale ratios. Asymptotic analysis then constructs limiting models and boundary layers, locating the validity of the preceding model families and guiding numerical resolution and error analysis.
\end{chaptersummary}
\learninggoals{Construct dimensionless parameters from characteristic scales; determine inner and outer regions by dominant balance; identify the fast scales removed by quasineutral, strong-field, and low-frequency reductions.}
\studyroute{The construction of matching conditions and proofs of uniform error estimates are omitted.}{singular perturbation; dominant balance; matched asymptotic expansions}{For $\epsilon u''+u'=0$ with $u(0)=0,u(1)=1$, construct inner and outer solutions, identify the boundary-layer location, and obtain a uniform approximation.}

\section{Nondimensionalization, Scale Ordering, and Asymptotic Expansions}\index{Singular perturbation}
Common dimensionless groups include
\[
\mathrm{Kn}=\frac{\lambda}{L},
\qquad
\mathrm{Re}=\frac{\rho UL}{\mu},
\qquad
\mathrm{Ma}=\frac{U}{c}.
\]
These parameters compare the magnitudes of different physical mechanisms. Model selection determines which terms remain in a specified limit.

\begin{authornote}
\textbf{Potential Question:}
"Why discuss dimensionless groups immediately after continuum equations? How do these numbers determine model selection?"

\textbf{Author Supplement:}
Numerical magnitudes of dimensional quantities depend on the chosen units and cannot be compared directly. After characteristic length, velocity, and time scales are selected, nondimensionalization expresses relative term sizes through dimensionless parameters. Their magnitudes indicate potentially dominant mechanisms and limits admitting reduced models. The scales themselves must be chosen for the particular problem.

\textbf{Core Conceptual Connection:}
\[
\text{continuum equations}
\longrightarrow
\text{dimensionless ratios}
\longrightarrow
\text{model regimes}
\]
\end{authornote}

Taylor expansion describes local variation through derivatives; asymptotic expansion orders scales in a small-parameter limit. Begin with the Taylor form:
\[
f(x+h)
=
f(x)+hf'(x)+\frac{h^2}{2}f''(x)+\cdots.
\]

\begin{askai}
Ask AI:
\begin{quote}
Compare Taylor and asymptotic expansions. Give an example of a useful divergent asymptotic series and explain why a smaller remainder at each fixed truncation order can be more relevant than convergence of the infinite sum.
\end{quote}

Follow-up:
\begin{quote}
For a simple function, calculate truncation errors at different orders and illustrate optimal truncation near the smallest term.
\end{quote}
\end{askai}

Taylor expansion is particularly useful for analyzing local differences between discrete and continuous operators.

An asymptotic expansion is more general. As \(\epsilon\to0\), suppose the scale functions satisfy
\[
\phi_{n+1}(\epsilon)=o(\phi_n(\epsilon)),
\]
and, for every fixed \(N\),
\[
f(\epsilon)
-
\sum_{n=0}^{N}a_n\phi_n(\epsilon)
=
o(\phi_N(\epsilon)),
\]
Then write
\[
f(\epsilon)\sim\sum_{n=0}^{\infty}a_n\phi_n(\epsilon).
\]

At each fixed truncation order, an asymptotic expansion requires a remainder smaller than the last retained scale. The series itself may diverge, and its scales may include fractional powers, logarithms, and exponentials:
\[
1,\quad
\epsilon^{1/2},\quad
\epsilon\log\epsilon,\quad
\ee^{-1/\epsilon}.
\]

A Taylor series is thus a special expansion generated by local derivatives in ordinary powers, whereas an asymptotic expansion emphasizes an ordering of scales.

Rapidly growing coefficients can cause an asymptotic series to improve initially and then deteriorate. For fixed \(\epsilon\), the terms
\[
|a_n\epsilon^n|
\]
may first decrease and then increase. Truncate near the smallest term rather than pursuing an infinite sum; this is optimal truncation.

If substituting
\[
u=u_0+\epsilon u_1+\epsilon^2u_2+\cdots
\]
preserves the differential order and boundary-condition structure at successive orders, the perturbation is regular.

If a small parameter multiplies the highest derivative, for example,
\[
\epsilon u''+u'=0,
\]
the limit \(\epsilon\to0\) lowers the order, and the outer solution cannot satisfy all boundary conditions. This is a singular perturbation, typically requiring boundary layers, matched asymptotic expansions, multiple scales, or WKB methods.

\section{Plasma Example: Small Debye Length and the Quasineutral Limit}

Consider the nondimensional one-dimensional Poisson equation
\[
-\epsilon^2\frac{\partial^2\phi}{\partial x^2}
=
n_i-n_e,
\qquad
\epsilon=\frac{\lambda_D}{L}\ll1.
\]
Here \(\lambda_D\) is the Debye length and \(L\) the macroscopic scale.

Setting \(\epsilon=0\) directly gives
\[
n_i=n_e,
\]
the quasineutrality condition. Net charge density is a higher-order quantity in the bulk.

If the boundary potential is incompatible with the bulk quasineutral solution, the small coefficient of the highest derivative makes the problem singular. An outer expansion in the bulk is
\[
\phi^{\mathrm{out}}
=
\phi_0+\epsilon\phi_1+\epsilon^2\phi_2+\cdots,
\]
For a layer with $O(1)$ potential variation and normalized net charge, the balance $\epsilon^2/\delta^2=O(1)$ gives $\delta=O(\epsilon)$. Under this ordering, introduce a stretched coordinate near the boundary:
\[
\xi=\frac{x}{\epsilon}.
\]

\begin{authornote}
\textbf{Potential Question:}
"Why introduce \(\xi=x/\epsilon\) near the boundary rather than an arbitrary new coordinate?"

\textbf{Author Supplement:}
A small coefficient can make the highest-derivative term negligible in the bulk, but the reduced equation may not satisfy all boundary conditions. Near a boundary, rapid variation over a short distance restores that derivative term to the same order as the others. Estimate all terms using an unknown layer thickness, determine it by dominant balance, and only then introduce the stretched coordinate.

\textbf{Core Conceptual Connection:}
\[
\text{singular perturbation}
\longrightarrow
\text{boundary-layer scale}
\longrightarrow
\xi=x/\epsilon
\]
\end{authornote}

At this scale,
\[
\frac{\partial}{\partial x}
=
\frac{1}{\epsilon}\frac{\partial}{\partial\xi},
\]

\begin{askai}
Ask AI:
\begin{quote}
Introduce an unknown boundary-layer width \(\delta(\epsilon)\) and determine how \(\delta\) scales with \(\epsilon\) by dominant balance. Then compare the stretched coordinate \(x/\epsilon\) with that scale and explain \(\xi=x/\epsilon\).
\end{quote}

Follow-up:
\begin{quote}
Draw a conceptual diagram showing the outer solution, inner solution, and overlap region, and identify where quasineutrality is valid.
\end{quote}
\end{askai}

The Poisson equation becomes
\[
-\frac{\partial^2\phi^{\mathrm{in}}}{\partial\xi^2}
=
n_i-n_e.
\]

Under this ordering, the Debye-scale charge-separation layer matches the wall potential to the bulk solution. Its profile also requires particle response, fluxes, and wall conditions. Other potential amplitudes or transport orderings may introduce additional transition layers.

This example illustrates characteristic features of asymptotic analysis:
\begin{enumerate}
  \item The small parameter is an explicit scale ratio, \(\lambda_D/L\).
  \item Bulk and boundary-layer regions use different coordinates.
  \item Setting the small parameter to zero lowers the differential order.
  \item Inner and outer solutions must match in an overlap region.
  \item The limiting model depends on the spatial region and boundary conditions.
\end{enumerate}

Related structures occur in low-frequency and strong-field plasma limits, drift kinetics, gyrokinetics, and MHD. Identify the small parameter, variables ordered at the same magnitude, any reduction in differential order, and any fast-time, short-space, or boundary-layer scales.

Common asymptotic notation has the following meanings:
\[
f=g+\Order(\epsilon^p)
\]
means that the remainder is bounded by a constant times \(\epsilon^p\);
\[
f=g+o(\epsilon^p)
\]
means that the remainder divided by \(\epsilon^p\) tends to zero;
\[
f\sim g
\]
may mean \(f/g\to1\) or an order-by-order asymptotic expansion, depending on context.

\section{Strong Fields, Low Frequencies, and Gyroscales}
A charged particle's gyrofrequency and characteristic gyroradius are
\[
\Omega_s=\frac{|q_s|B_0}{m_s},\qquad
\rho_{Ls}=\frac{v_{\mathrm{th},s}}{\Omega_s}.
\]
When $\rho_{Ls}/L\ll1$ and $\omega/\Omega_s\ll1$, variables can be organized into fast gyrophase and slow guiding-center motion.
Drift kinetics further assumes $k_\perp\rho_{Ls}\ll1$; gyrokinetics permits perturbations with $k_\perp\rho_{Ls}=O(1)$,
retaining finite-Larmor-radius effects. Chapter~\ref{ch:19} develops the five-dimensional distribution, gyroaveraging, and self-consistent field equations.

\part{Mathematical Structure of Continuum Equations}
\chapter{Function Spaces, Sobolev Spaces, and Weak Solutions}\label{ch:5}
\begin{chaptersummary}
The first four chapters established physical models and their scale assumptions. Solving these equations also requires specifying the meaning of derivatives, boundary values, and solutions. Starting from integration by parts for Poisson's equation, this chapter introduces weak derivatives and Sobolev spaces, uses energy estimates to establish existence and stability, and discusses the selection of discontinuous solutions of conservation laws. The next chapter constructs finite-dimensional residual constraints in these spaces.
\end{chaptersummary}
\learninggoals{Derive a weak form including boundary terms; define weak derivatives and traces; apply the Poincar\'e inequality and Lax--Milgram theorem under explicit assumptions; distinguish broken from global spaces.}
\studyroute{Weak forms, space definitions, and basic estimates are given; proofs of general trace theorems, compact embeddings, and low-regularity elliptic theory are omitted.}{weak derivative; Sobolev trace; Poincar\'e inequality; Lax--Milgram theorem}{Integrate the Poisson problem on the unit interval by parts, then use an energy estimate to show how $H^{-1}$ data control the $H_0^1$ solution.}

\section{From Strong Form to an Integral Identity}\index{Weak form}
Let $\Omega\subset\R^d$ be a bounded Lipschitz domain. First consider a sufficiently smooth solution of
\[
-\Delta u=f\quad\text{in }\Omega,\qquad u=0\quad\text{on }\partial\Omega.
\]
Multiply by a test function $v$ and integrate. The divergence theorem gives
\begin{equation}\label{eq:poisson-green}
\int_\Omega\nabla u\cdot\nabla v\dd x
=\int_\Omega fv\dd x+\int_{\partial\Omega}(\partial_nu)v\dd S.
\end{equation}
For test functions with zero boundary values, the boundary term vanishes. This identity uses only first derivatives of $u$,
so it can be interpreted in a larger space than the classically twice-differentiable functions.
Nonhomogeneous Dirichlet data are incorporated through a prescribed lifting $u_g$, writing $u=u_g+w$ while retaining the zero-trace test space.
Neumann data enter the boundary pairing on the right.

\conceptbridge{Why specify a function space after reducing the derivative order?}
{Integration by parts changes the regularity required by the equation. Square-integrable first weak derivatives make the energy integral meaningful; traces specify boundary conditions. The function space determines admissible solutions, testing directions, and the norm used in stability estimates.}
{Integration by parts $\longrightarrow$ weak derivatives and boundary traces $\longrightarrow$ a variational problem.}

\section{Weak Derivatives and Completeness}\index{Sobolev spaces}
For locally integrable $u$, if
\[
\int_\Omega u\,\partial_i\varphi\dd x=-\int_\Omega g_i\varphi\dd x
\quad\text{for every }\varphi\in C_c^\infty(\Omega),
\]
then $g_i$ is the $i$th weak derivative of $u$. A distributional derivative may be a general distribution;
for example, the derivative of a one-dimensional step function is a Dirac distribution, which is not in $L^2$.

A normed space is a Banach space if every Cauchy sequence converges within it.
A complete space whose norm is induced by an inner product is a Hilbert space. Completeness ensures that limits obtained from energy estimates remain admissible functions.
\[
L^2(\Omega)=\{u:\int_\Omega|u|^2\dd x<\infty\},\qquad
(u,v)_{L^2}=\int_\Omega uv\dd x.
\]
Functions equal almost everywhere are identified. Define
\[
H^1(\Omega)=\{u\in L^2(\Omega):\partial_i u\in L^2(\Omega),\ 1\le i\le d\},
\qquad
\norm{u}_{H^1}^2=\norm{u}_{L^2}^2+\norm{\nabla u}_{L^2}^2.
\]
For integers $m\ge0$, $H^m$ requires all weak derivatives with $|\alpha|\le m$ to belong to $L^2$.
Real-order spaces on $\R^d$ can be defined through the Fourier transform:
\[
\norm{u}_{H^s(\R^d)}^2
=\int_{\R^d}(1+|\bm k|^2)^s|\widehat u(\bm k)|^2\dd\bm k,
\]
Here the Fourier transform is unitarily normalized; $H^s$ on a domain can be constructed using restriction and extension.

\section{Traces, the Poincar\'e Inequality, and Energy}\index{Trace operator}
On the stated domain, the trace operator $\gamma:H^1(\Omega)\to H^{1/2}(\partial\Omega)$
continuously extends the boundary restriction of smooth functions. Define
\[
H_0^1(\Omega)=\overline{C_c^\infty(\Omega)}^{H^1}
=\{u\in H^1(\Omega):\gamma u=0\}.
\]
For $v\in H_0^1(\Omega)$, the Poincar\'e inequality gives
\[
\norm{v}_{L^2}\le C_P\norm{\nabla v}_{L^2}.
\]
Thus $\norm{\nabla v}_{L^2}$ is an energy norm equivalent to the $H^1$ norm on this space.
A pure Neumann problem retains a constant nullspace, requires fixing the mean, and must satisfy
$\int_\Omega f\dd x+\int_{\partial\Omega}g_N\dd S=0$.

\section{The Weak Problem and the Lax--Milgram Theorem}
Let $V=H_0^1(\Omega)$ and $V'=H^{-1}(\Omega)$. The weak Poisson problem is
\[
\text{Find }u\in V:\qquad b(u,v)=\ell(v)\quad\forall v\in V,\qquad
b(u,v)=\int_\Omega\nabla u\cdot\nabla v\dd x.
\]
For $f\in L^2$, $\ell(v)=\int fv$; general $f\in H^{-1}$ requires a duality pairing.
\begin{theorem}[Lax--Milgram]
Let $V$ be a real Hilbert space and let a bilinear form satisfy
$|b(w,v)|\le M\norm{w}_V\norm{v}_V$ and
$b(v,v)\ge\alpha\norm{v}_V^2$, where $\alpha>0$.
If $\ell\in V'$, there is a unique $u\in V$ satisfying $b(u,v)=\ell(v)$, with
$\norm{u}_V\le\norm{\ell}_{V'}/\alpha$.
\end{theorem}
For the Poisson problem measured in the energy norm, $M=\alpha=1$.
Testing with $v=u$ gives this stability estimate; existence follows from Riesz representation and a bounded coercive operator.
Convection-dominated, saddle-point, and wave problems require estimates adapted to their operators.

\aiexercise{On $(0,1)$, take $u(x)=x(1-x)$. Derive the weak form from $-u''=2$ and calculate the energy and right-hand-side pairing. Then verify that a continuous piecewise-linear function with a corner belongs to $H_0^1$.}
{Compute the distributional derivative of a one-dimensional step function and explain why it is not in $H^1$. Next derive the compatibility condition and constant nullspace of a pure Neumann problem.}

\section{Weak Solutions and Discontinuities in Conservation Laws}
For $u_t+\partial_xF(u)=0$ in one dimension, a test function
$\varphi\in C_c^\infty([0,T)\times\R)$ gives
\[
\int_0^T\int_\R(u\varphi_t+F(u)\varphi_x)\dd x\dd t
+\int_\R u_0(x)\varphi(0,x)\dd x=0.
\]
A moving discontinuity $x=s t$ satisfies the Rankine--Hugoniot condition
\[
s(u_R-u_L)=F(u_R)-F(u_L).
\]
The weak identity expresses conservation across a discontinuity. For a one-dimensional scalar conservation law with suitable data, a family of convex entropy inequalities selects the unique entropy solution.
An entropy pair $\eta,q$ satisfies $q'=\eta'F'$ and requires, in distributions,
$\partial_t\eta(u)+\partial_xq(u)\le0$. The finite-volume fluxes of Chapter~\ref{ch:8} approximate this propagation.

\section{Broken Sobolev Spaces}
For a mesh $\mathcal T_h$, define
\[
H^1(\mathcal T_h)=\{v\in L^2(\Omega):v|_K\in H^1(K),\ K\in\mathcal T_h\}.
\]
Derivatives are initially evaluated elementwise, and interface traces may jump. Distributional derivatives also contain interface contributions;
DG numerical fluxes or elliptic penalty terms account for these interactions.
A broken space specifies admissible discontinuities; the algorithm must additionally define residuals and interface coupling.

\chapter{Projection, Interpolation, and the Galerkin Framework}\label{ch:6}
\begin{chaptersummary}
Weak formulations specify continuous solution and test spaces. Discretization then selects a finite-dimensional representation and residual constraints. Comparing interpolation, orthogonal projection, Galerkin testing, collocation, and control-volume conservation exposes these choices, while error orthogonality links approximation power to solution accuracy. Subsequent methods define their unknowns, operators, and interfaces within this framework.
\end{chaptersummary}
\learninggoals{Specify trial and test spaces; distinguish five approximation or residual mechanisms; state the conditions for an energy projection; construct mass matrices, residuals, and sparse coupling from local bases.}
\studyroute{Basic error estimates for coercive problems are given; general discrete inf--sup theory for Petrov--Galerkin methods and nonconforming error estimates are omitted.}{Galerkin orthogonality; Riesz representation; Cea lemma; Petrov--Galerkin inf-sup}{Compare energy errors for $-\!u''=f$ on two linear finite-element meshes, then use a convection problem to examine stability when the test space changes.}

\section{From a Continuous Operator to Residual Constraints}
Let the continuous problem be $\mathcal A u=f$, with admissible fields in $V$ and test space $W$.
Its weak form is $b(u,v)=\ell(v)$.
Choose finite-dimensional trial and test spaces $V_h$ and $W_h$:
\begin{equation}\label{eq:galerkin-general}
\text{Find }u_h\in V_h:\qquad b(u_h,v_h)=\ell(v_h)
\quad\forall v_h\in W_h.
\end{equation}
The residual functional $r_h(v)=\ell(v)-b(u_h,v)$ vanishes on $W_h$;
it generally remains nonzero in other testing directions. Equal trial and test spaces give a Galerkin method;
different spaces give a Petrov--Galerkin method. For nonlinear problems, the residual depends nonlinearly on the solution.

\conceptbridge{After choosing a finite-dimensional representation, what is still needed to obtain numerical equations?}
{Basis functions specify candidate solutions. Numerical equations additionally require the residual to vanish in selected testing directions, at selected points, or over selected control volumes. This constraint mechanism, together with interface treatment, defines the discrete operator.}
{Representation space $\longrightarrow$ residual constraints $\longrightarrow$ finite-dimensional equations.}

\section{Interpolation and Orthogonal Projection}\index{Galerkin method}\index{Interpolation}
An interpolation operator $I_h$ satisfies prescribed degrees of freedom $\lambda_i$:
$\lambda_i(I_hu)=\lambda_i(u)$. Nodal interpolation requires well-defined point values;
point evaluation need not be continuous on a general higher-dimensional $H^1$ space, where local-average degrees of freedom may be used instead.

An orthogonal projection $P_h$ in a Hilbert space satisfies
\[
P_hu\in V_h,\qquad (u-P_hu,v_h)_V=0\quad\forall v_h\in V_h.
\]
For $u(x)=x^2$ and $V_h=\Poly_1(-1,1)$, endpoint interpolation gives $I_hu=1$,
whereas $L^2$ orthogonal projection gives $P_hu=1/3$.
The same representation space can support different approximation operators.

\section{Riesz Representation, Energy Norms, and Error Estimates}
The Riesz map $R_V:V\to V'$ of a Hilbert space is defined by
$\langle R_Vw,v\rangle=(w,v)_V$,
giving each continuous linear functional a unique inner-product representative. A PDE bilinear form must additionally be checked for symmetry and coercivity.
For a conforming Galerkin approximation with $V_h\subset V$, subtracting the continuous and discrete equations gives
\[
b(u-u_h,v_h)=0\quad\forall v_h\in V_h.
\]
If $b$ is symmetric positive definite, it defines the energy norm $\norm{v}_b^2=b(v,v)$, and
\[
\norm{u-u_h}_b=\inf_{w_h\in V_h}\norm{u-w_h}_b.
\]
For a bounded, coercive, but not necessarily symmetric $b$, any $w_h\in V_h$ gives
\[
\alpha\norm{u-u_h}_V^2
\le b(u-u_h,u-w_h)
\le M\norm{u-u_h}_V\norm{u-w_h}_V.
\]
Canceling the error norm yields C\'ea's estimate:
\[
\norm{u-u_h}_V\le\frac{M}{\alpha}
\inf_{w_h\in V_h}\norm{u-w_h}_V.
\]
This relates numerical error to the approximation power of the space. Petrov--Galerkin and noncoercive problems require discrete inf--sup conditions.

\section{Degrees of Freedom and Residual Mechanisms}
\begin{center}\small
\begin{tabularx}{\textwidth}{@{}l>{\raggedright\arraybackslash}X>{\raggedright\arraybackslash}X@{}}
\toprule
Method & Representation and degrees of freedom & Equation or residual constraint\\
\midrule
FDM & Grid functions, nodal values & Difference equations at grid points\\
FVM & Cell averages and local reconstruction & Integral conservation in each control volume\\
FEM & Local bases, trace-compatible degrees of freedom & Weak form on a finite test space\\
Spectral Galerkin & Global modal coefficients & Residual vanishes against retained global modes\\
Spectral collocation & Nodal values of a global interpolant & Differential equation at collocation points\\
DG / FR & Element-local polynomials & Weak residual / corrected flux derivatives and shared interface fluxes\\
LBM & Populations at discrete velocities & Streaming, collision, and their conserved moments\\
\bottomrule
\end{tabularx}
\end{center}
These methods share the language of finite-dimensional representations, residuals, and structural analysis.
An interpretation as Hilbert-space orthogonal projection depends on the inner product, trial space, and residual conditions.

\aiexercise{Compute endpoint interpolation and $L^2$ projection of $x^2$ onto $\Poly_1$ on $[-1,1]$, including both errors. Identify the inner product underlying orthogonality in symmetric Poisson Galerkin approximation.}
{Compare central differences, cell-average FVM, and linear FEM for the same Poisson problem. Identify their degrees of freedom and residual constraints, and state the quadrature and boundary assumptions needed to obtain identical matrices on a uniform grid.}

\section{Formation of Finite-Dimensional Algebraic Systems}
Discretization of a PDE typically gives
\[
\bm A\bm x=\bm b
\]
or a nonlinear residual system
\[
\bm R(\bm U)=0.
\]

\begin{authornote}
\textbf{Potential Question:}
"Why does a continuous PDE become a finite-dimensional vector equation? Where do the degrees of freedom, residuals, and matrix entries come from?"

\textbf{Author Supplement:}
A continuous field has a value at every point, whereas a computer stores finitely many parameters. Discretization first represents an approximate field by grid values or basis coefficients, then enforces the equation in finitely many senses: collocation, control-volume conservation, or integration against test functions. The result is a finite-dimensional algebraic problem; nonlinear or time-dependent problems additionally require linearization or time discretization.

\textbf{Core Conceptual Connection:}
\[
\text{continuous PDE}
\longrightarrow
\text{finite-dimensional DOFs}
\longrightarrow
\text{algebraic system}
\]
\end{authornote}

Local differential operators usually produce sparse matrices, whereas boundary integral operators often produce dense ones. Matrix structure reflects the range of physical coupling.

\begin{authornote}
\textbf{Potential Question:}
"Why does differential-operator locality become matrix sparsity? Why are integral-equation matrices often dense?"

\textbf{Author Supplement:}
A matrix row records which degrees of freedom enter one discrete equation. A local differential operator with a local basis or finite-difference stencil usually couples only nearby degrees of freedom, leaving few nonzeros per row. Global bases or integral operators may couple one degree of freedom to many distant ones. Sparsity therefore depends on both the operator and the discretization.

\textbf{Core Conceptual Connection:}
\[
\text{operator locality}
\longrightarrow
\text{matrix sparsity or density}
\]
\end{authornote}


With basis functions $\phi_j$ and $u_h=\sum_jU_j\phi_j$, equation~\eqref{eq:galerkin-general} gives
$A_{ij}=b(\phi_j,\psi_i)$ and $b_i=\ell(\psi_i)$.
For time-dependent problems,
\[
M_{ij}=(\phi_j,\psi_i),\qquad M_h\dot U_h=R_h(U_h,t).
\]
The mass matrix represents the pairing of the time derivative; the residual contains volume, boundary, and source terms.
Their structure follows from representation and constraints. Algebraic solution algorithms are developed in Volume~0.

\chapter{Finite Differences: Consistency, Stability, and Wave Errors}\label{ch:7}
\begin{chaptersummary}
With function spaces and residual mechanisms established, grid functions and difference operators provide a direct discrete model. Taylor expansion gives consistency; error evolution gives stability; modified equations and numerical frequencies quantify dissipation and dispersion. These analyses provide a common reference for finite-volume and high-order methods and time integration.
\end{chaptersummary}
\learninggoals{Derive truncation errors from stencils; state the assumptions of periodic Fourier analysis; distinguish stability, the necessary CFL condition, and convergence; calculate amplitude and phase errors.}
\studyroute{Advection, grid norms, and error propagation establish the difference analysis. Finite-boundary SBP--SAT analysis and the full convergence proof are left for further study.}{grid function; truncation error; von Neumann analysis; modified equation; SBP--SAT}{Derive the upwind amplification factor mode by mode, then use discrete Parseval to prove periodic discrete L2 contraction.}
\section{Grid Functions and Difference Operators}
On uniform nodes $x_i=ih$, let $U_i\approx u(x_i)$. Define
\[
(D_-U)_i=\frac{U_i-U_{i-1}}h,\qquad
(D_0U)_i=\frac{U_{i+1}-U_{i-1}}{2h},\qquad
\norm{U}_{2,h}^2=h\sum_i|U_i|^2.
\]
Grid functions belong to finite-dimensional sequence spaces. Difference operators approximate the differential equation at grid points;
interpolation can compare a grid function with a continuous field.
\section{Consistency, Stability, and Modified Equations}\index{Finite difference method}
Consider linear advection:
\[
u_t+au_x=0,
\qquad
u(x,t)=u_0(x-at).
\]

\begin{authornote}
\textbf{Potential Question:}
"After discussing complex physical models, why begin numerical discretization with linear advection?"

\textbf{Author Supplement:}
Complex transport equations combine several physical mechanisms, obscuring the source of numerical error. Linear advection isolates propagation at a fixed speed, allowing separate examination of changes in transport, amplitude, and phase. It is a reference model for numerical mechanisms; conclusions must be reassessed for nonlinearity, variable coefficients, and boundaries.

\textbf{Core Conceptual Connection:}
\[
\text{complicated transport physics}
\longrightarrow
\text{linear advection as minimal model}
\]
\end{authornote}

The forward time difference satisfies
\[
\frac{u_i^{n+1}-u_i^n}{\Delta t}
=
u_t+\Order(\Delta t),
\]
The centered spatial difference satisfies
\[
\frac{u_{i+1}-u_{i-1}}{2\Delta x}
=
u_x+\Order(\Delta x^2).
\]
The order of a scheme follows from matching local expansions of discrete and continuous operators.

For a linear constant-coefficient scheme, write an error mode as
\[
e_i^n=\widehat e^n\ee^{\ii kx_i},
\]

\begin{authornote}
\textbf{Potential Question:}
"Why can arbitrary errors be analyzed through Fourier modes? Why does studying one mode characterize overall error propagation?"

\textbf{Author Supplement:}
For a linear, constant-coefficient, translation-invariant scheme, each discrete Fourier mode retains its wave number after one update; only its complex amplitude changes. Amplification factors can therefore be studied separately to determine modal growth. This analysis applies most directly to infinite or periodic uniform grids. Finite boundaries, variable coefficients, and nonlinearity can introduce additional coupling.

\textbf{Core Conceptual Connection:}
\[
\text{arbitrary numerical error}
\longrightarrow
\text{Fourier modes}
\longrightarrow
\text{amplification factor}
\]
\end{authornote}

\begin{askai}
Ask AI:
\begin{quote}
Starting from the Fourier decomposition of an arbitrary discrete error, develop von Neumann stability analysis step by step: explain why a linear constant-coefficient difference scheme preserves each Fourier mode and how this permits separate amplification-factor analysis.
\end{quote}

Follow-up:
\begin{quote}
Substitute \(e_i^n=\widehat e^n e^{\ii kx_i}\) into first-order upwinding, derive \(G(k)\), and plot \(|G(k)|\) for different CFL numbers.
\end{quote}
\end{askai}

This gives
\[
\widehat e^{n+1}=G(k)\widehat e^n.
\]
The one-step amplification factor describes each mode. For an advection discretization required to be non-growing, the contraction condition is
\[
|G(k)|\le1.
\]

Explicit schemes are also subject to a CFL constraint. Define
\[
\nu=\frac{a\Delta t}{\Delta x}.
\]
The numerical domain of dependence must contain the physical domain of dependence.

\begin{authornote}
\textbf{Potential Question:}
"How is \(\nu=a\Delta t/\Delta x\) related to domains of dependence, and why does it constrain explicit schemes?"

\textbf{Author Supplement:}
Characteristic speeds determine physical dependence, while the stencil and time step determine numerical information propagation. Containment of the physical domain by the numerical domain under refinement is necessary for convergence; error-evolution estimates establish stability. CFL connects these two propagation ranges.

\textbf{Core Conceptual Connection:}
\[
\text{amplification factor stability}
\longrightarrow
\text{CFL number}
\longrightarrow
\text{domain of dependence}
\]
\end{authornote}

For \(a>0\), first-order upwinding is
\[
u_i^{n+1}
=
u_i^n-\nu(u_i^n-u_{i-1}^n)
=
(1-\nu)u_i^n+\nu u_{i-1}^n.
\]
When
\[
0\le\nu\le1
\]
the new value is a convex combination of old values, so no new local extrema are created.

Its modified equation is
\[
u_t+au_x
=
\frac{a\Delta x}{2}(1-\nu)u_{xx}
+\cdots.
\]

\begin{askai}
Ask AI:
\begin{quote}
Perform modified-equation analysis for first-order upwinding. Starting from Taylor expansion, derive term by term
\[
u_t+au_x=\frac{a\Delta x}{2}(1-\nu)u_{xx}+\cdots
\]
and explain how the diffusion term on the right represents numerical dissipation.
\end{quote}

Follow-up:
\begin{quote}
Write a minimal Python experiment comparing dissipation in centered and upwind schemes for a propagating Gaussian pulse.
\end{quote}
\end{askai}

The equivalent diffusion explains upwind dissipation: it suppresses high-frequency error and also smooths genuine small scales.
General linear stability requires, over a fixed physical time $T$,
$\sup_{n\Delta t\le T}\norm{S_h^n}\le C_T$, with a mesh-independent constant.
For scalar translation-invariant updates, Parseval converts this to a uniform bound on modal powers.
Bounded growth with $|G|\le1+C\Delta t$ is admissible; a one-step value $|G|>1$ must be assessed in the refinement limit.
For example,
\[
|G(k)|>1,
\]
indicates one-step growth. If an amplification value remains separated from 1 by a fixed positive amount under refinement, its powers become unbounded over fixed physical time.
For a well-posed linear initial-value problem in a compatible discrete setting, consistency and stability together imply convergence.

\section{Numerical Dispersion and Wave Errors}
Let a continuous wave be
\[
u=\widehat u\ee^{\ii(kx-\omega t)}.
\]
For linear advection,
\[
\omega=ak.
\]
the phase and group velocities both equal \(a\):
\[
c_p=\frac{\omega}{k}=a,
\qquad
c_g=\frac{\dd\omega}{\dd k}=a.
\]

A discrete scheme generally gives a numerical frequency
\[
\omega_{\mathrm{num}}(k).
\]

\begin{askai}
Ask AI:
\begin{quote}
Choose a specific difference scheme for one-dimensional advection and derive \(\omega_{\mathrm{num}}(k)\). Explain how \(\operatorname{Re}\omega_{\mathrm{num}}\) and \(\operatorname{Im}\omega_{\mathrm{num}}\) describe phase and dissipation errors.
\end{quote}

Follow-up:
\begin{quote}
Compare phase-speed errors for two schemes at the same wave number \(k\), and estimate the accumulated phase shift after propagation over many wavelengths.
\end{quote}
\end{askai}

If
\[
\operatorname{Re}\omega_{\mathrm{num}}\neq ak,
\]
there is numerical dispersion; if
\[
\operatorname{Im}\omega_{\mathrm{num}}<0,
\]
there is numerical dissipation.

Shock calculations generally require some dissipation, whereas acoustic, electromagnetic, and elastic waves place greater demands on long-range phase accuracy. Scheme selection must reflect the target physics.

\chapter{Finite Volumes: Numerical Fluxes and High-Resolution Schemes}\label{ch:8}
\begin{chaptersummary}
The weak form of a conservation law describes discontinuities through control-volume flux balances. Cell averages, Riemann problems, and shared fluxes lead to finite-volume methods; Godunov's order barrier motivates nonlinear reconstruction. Conservation, monotonicity, and entropy stability impose distinct constraints on fluxes and updates, providing the basis for DG and FR interface design.
\end{chaptersummary}
\learninggoals{Construct cell updates from integral conservation; explain Riemann fluxes; distinguish Godunov's theorem, TVD, and high-order reconstruction; check reconstructed averages and units.}
\studyroute{The chapter develops Godunov, MUSCL, TVD, and WENO in sequence. Characteristic decomposition for systems, full smoothness indicators, and nonlinear entropy-convergence proofs are left for further study.}{finite volume; Riemann problem; Godunov order barrier; MUSCL; WENO}{Calculate a minmod reconstruction from three cell averages, then apply the Rusanov flux to two adjacent control volumes and check the signs of their internal fluxes.}
\section{Control Volumes, Cell Averages, and Shared Fluxes}\index{Finite volume method}
Finite-volume methods start from control-volume conservation:
\[
\frac{\dd}{\dd t}
\int_{\Omega_i}U\dd V
+
\int_{\partial\Omega_i}
\bm F\cdot\bm n\dd S
=
\int_{\Omega_i}S\dd V.
\]
Adjacent cells use the same interface flux with opposite signs. Internal fluxes therefore cancel in a global sum, giving local and global conservation.


Define $\bar U_K=|K|^{-1}\int_KU$. Degrees of freedom are cell averages; interface states are reconstructed from them.
A consistent numerical normal flux satisfies $\widehat F(U,U;n)=F(U)\cdot n$ and changes sign when the states are exchanged and the normal is reversed.
For a one-dimensional scalar equation, Riemann initial data are constant states $U_L,U_R$ on either side of an interface.
The exact interface flux is the Godunov flux; approximate solvers use computable wave-speed estimates.
For example,
\[
\widehat F(U_L,U_R)=\frac{F(U_L)+F(U_R)}2-\frac{\alpha}2(U_R-U_L),
\qquad
\alpha\ge\max_{[U_L,U_R]}|F'(U)|.
\]
For linear advection $F(U)=aU$ with $a>0$, the exact flux is $aU_L$.
A cell update also integrates the flux over one time step. High-order spatial reconstruction and time integration jointly determine fully discrete accuracy.

\section{Godunov's Theorem and Nonlinear Reconstruction}
Godunov's theorem is a basic limitation in numerical conservation laws. For linear scalar advection, consider a scheme that is

\begin{enumerate}
  \item linear;
  \item monotone, creating no new extrema;
  \item consistent.
\end{enumerate}

For general noninteger Courant numbers, a linear monotone consistent scheme has at most first-order accuracy. Exact grid translation at integer Courant numbers for constant-speed advection is a special case.

Higher-order linear reconstruction generally requires alternating positive and negative weights, whereas monotonicity requires a convex-combination update. These requirements are generally incompatible.

Improving the accuracy of monotone schemes therefore requires changing these conditions. Modern high-resolution methods commonly use nonlinear reconstruction:

\begin{authornote}
\textbf{Potential Question:}
"Why does Godunov's theorem motivate nonlinear reconstruction? Which restriction does nonlinearity remove?"

\textbf{Author Supplement:}
Under the theorem's assumptions, linear monotone schemes are at most first order. High-resolution methods therefore make reconstruction or limiting depend on the local solution: retain higher order in smooth regions and adjust slopes near steep gradients to suppress oscillation. The scheme becomes nonlinear and leaves the class of linear schemes restricted by the theorem.

\textbf{Core Conceptual Connection:}
\[
\text{Godunov theorem}
\longrightarrow
\text{nonlinear reconstruction}
\longrightarrow
\text{high-resolution schemes}
\]
\end{authornote}

\begin{itemize}
  \item Preserve conservation.
  \item Allow a nonlinear algorithm.
  \item Recover high order in smooth regions.
  \item Reduce to a first-order monotone scheme near discontinuities.
\end{itemize}

For a finite-volume method,
\[
\frac{\dd \bar U_i}{\dd t}
=
-\frac{1}{\Delta x}
\left(
\widehat F_{i+1/2}
-
\widehat F_{i-1/2}
\right).
\]
the key step is constructing left and right interface states
\[
U_{i+1/2}^{L},
\qquad
U_{i+1/2}^{R},
\]
and evaluating a numerical flux with an exact or approximate Riemann solver.

MUSCL uses piecewise-linear reconstruction, for example,
\[
U_{i+1/2}^{L}
=
U_i+\frac12\sigma_i,
\]

\begin{askai}
Ask AI:
\begin{quote}
Construct MUSCL reconstruction from three neighboring cell averages. Calculate \(U_{i+1/2}^L\) using an unlimited slope and a minmod limiter, and explain how limiting retains second order in smooth regions and reverts to first order near discontinuities.
\end{quote}

Follow-up:
\begin{quote}
For data containing a discontinuity, calculate and compare minmod, van Leer, and Superbee slopes.
\end{quote}
\end{askai}

Here $\sigma_i$ is the increment across one cell width and has the same units as $U_i$.
For example, $\sigma_i=\operatorname{minmod}(U_i-U_{i-1},U_{i+1}-U_i)$,
where minmod returns the argument of smaller magnitude for equal signs, and zero otherwise.
The slope \(\sigma_i\) adapts to local smoothness; minmod, van Leer, and Superbee impose different restrictions. This retains approximately second-order behavior in smooth regions and reduces reconstruction order near shocks or discontinuities to suppress oscillations.

A TVD method requires nonincreasing discrete total variation:
\[
TV(U^{n+1})\le TV(U^n),
\qquad
TV(U)=\sum_i|U_{i+1}-U_i|.
\]
This turns suppression of spurious oscillations into an analyzable discrete condition.

ENO/WENO selects or weights several candidate stencils rather than directly limiting one slope. WENO uses a high-order combination in smooth regions and reduces weights on stencils containing a discontinuity.

Godunov's theorem exposes a structural tradeoff:
\[
\text{High-order accuracy}
\quad\leftrightarrow\quad
\text{Monotonicity and nonoscillation}.
\]
Scheme design adjusts reconstruction and dissipation to local solution smoothness.

\section{Conservation, Entropy, and High Order in Smooth Regions}
WENO combines candidate reconstructions $q_r$ with nonlinear weights:
\[
U_{i+1/2}^{L}=\sum_r\omega_rq_r,\qquad
\omega_r=\frac{\alpha_r}{\sum_s\alpha_s},\qquad
\alpha_r=\frac{d_r}{(\epsilon_W+\beta_r)^2}.
\]
$\beta_r$ measures stencil smoothness and $d_r$ gives the ideal weights of the corresponding high-order linear combination;
the scale of $\epsilon_W$ must be compatible with the variables and mesh.
Smooth-region order depends on how rapidly the weights recover their ideal values; extrema require separate analysis.
Fully discrete scalar TVD results depend on the flux, limiter, and time step.
Entropy-stable systems additionally require entropy variables and compatible numerical entropy fluxes. For the original WENO construction, see~\cite{jiang1996}.

\chapter{Finite Elements, Spectral Methods, and Integral Operators}\label{ch:9}
\begin{chaptersummary}
The Galerkin framework connects error analysis to space selection and testing. Local finite-element assembly leads to global Fourier, Legendre, and Chebyshev representations, followed by integral-operator discretization with MoM. Support determines coupling range, while regularity and quadrature govern error, linking approximation theory to matrix structure and solvers in Volume 0.
\end{chaptersummary}
\learninggoals{Assemble one-dimensional linear finite-element mass and stiffness matrices; distinguish spectral Galerkin, collocation, and pseudospectral methods; construct a MoM matrix from an integral kernel; explain how smoothness controls approximation rates.}
\studyroute{Higher-dimensional element families, general spectral-error proofs, singular integration, and full radiation-condition treatment are omitted. Matrix solvers and fast operator application are developed in Volume~0.}{Galerkin finite element; spectral collocation; Fourier pseudospectral; Chebyshev approximation; method of moments}{For one-dimensional Poisson, specify the degrees of freedom, boundary constraints, and matrix entries for local linear FEM and global Legendre Galerkin approximation.}
\section{Local Finite-Element Spaces and Assembly}\index{Finite element method}
Finite elements start from a weak form. For Poisson's equation,
\[
-\nabla^2u=f
\]
multiply by a test function and integrate by parts:
\[
\int_\Omega\nabla v\cdot\nabla u\dd V
=
\int_\Omega vf\dd V+\text{boundary terms}.
\]

\begin{askai}
Ask AI:
\begin{quote}
Discretize one-dimensional Poisson using second-order central differences and linear finite elements. Write both matrices for five interior nodes and compare how differential-operator locality becomes matrix sparsity.
\end{quote}

Follow-up:
\begin{quote}
Using the same Poisson problem, explain the relationships among strong form, weak form, test function, basis function, and stiffness matrix.
\end{quote}
\end{askai}

Expanding
\[
u_h=\sum_jU_j\phi_j
\]
gives
\[
\bm K\bm U=\bm F.
\]

The expansion turns the weak form into equations for basis coefficients. A one-dimensional linear element makes the local matrices explicit, after which shared nodes determine assembly.


On a one-dimensional element $[x_j,x_{j+1}]$ of length $h_j$, the linear shape functions are
$\phi_0=1-\xi$, $\phi_1=\xi$, $\xi=(x-x_j)/h_j$.
Direct integration gives
\[
K^{(j)}=\frac1{h_j}\begin{pmatrix}1&-1\\-1&1\end{pmatrix},
\qquad
M^{(j)}=\frac{h_j}6\begin{pmatrix}2&1\\1&2\end{pmatrix}.
\]
Shared nodes represent the same degree of freedom; local entries are added using the node map.
Conforming spaces satisfy the required trace conditions. Scalar $H^1$ elements serve Poisson problems;
Maxwell curl problems require the corresponding tangential-trace structure, developed in Volume~2.

\conceptbridge{How can the same global polynomial space produce both Galerkin and collocation methods?}
{Modal coefficients and nodal values are representation choices. Galerkin constrains residuals by integral pairing; collocation constrains differential residuals at selected nodes. Suitable quadrature may relate the resulting matrices, with the precise relationship determined by testing, integration, and boundary treatment.}
{Global approximation space $\longrightarrow$ modal or nodal degrees of freedom $\longrightarrow$ residual mechanism.}

\section{Global Modes, Collocation, and Pseudospectral Methods}\index{Spectral methods}
Local finite elements organize assembly through compact support; global spectral representations exploit overall smoothness for accuracy. On the periodic interval $[0,2\pi)$, Fourier Galerkin takes
\[
u_N(x)=\sum_{k=-N}^{N}\widehat u_k\ee^{\ii kx},\qquad
\widehat{(\partial_xu_N)}_k=\ii k\widehat u_k.
\]
Galerkin tests against the retained Fourier modes. Pseudospectral evaluation transforms modes to nodal values,
computes nonlinear products in physical space, and transforms back. Unresolved high wave numbers may alias into lower ones;
appropriate padding or truncation removes aliasing for quadratic nonlinearities.

On $[-1,1]$, Legendre polynomials satisfy
\[
\int_{-1}^1P_jP_k\dd x=\frac{2}{2k+1}\delta_{jk}.
\]
Chebyshev polynomials $T_k(\cos\theta)=\cos(k\theta)$
are orthogonal with weight $(1-x^2)^{-1/2}$. This weighted inner product must be distinguished from the unweighted physical $L^2$ inner product.
Spectral collocation can use Chebyshev--Lobatto nodes $x_j=\cos(j\pi/p)$,
with $D_{ij}=\ell_j'(x_i)$. Differential residuals are constrained at interior nodes, and boundary values at boundary nodes.

For sufficiently smooth solutions with suitable analytic continuation, increasing the global degree can yield rapid or exponential convergence.
Finite regularity or discontinuities generally yield algebraic convergence, with possible Gibbs oscillations near interfaces.
Local polynomials facilitate geometry adaptation and refinement; global modes exploit global smoothness.
Error comparisons should report total degrees of freedom, quadrature error, and temporal error.

\section{Integral Operators and the Method of Moments}\index{Method of moments}
The method of moments in computational electromagnetics typically starts from an integral equation:
\[
\Lop J=g
\]
Expand
\[
J=\sum_{n=1}^N I_nf_n,
\]

\begin{authornote}
\textbf{Potential Question:}
"Why switch from differential equations to an integral equation and a current expansion?"

\textbf{Author Supplement:}
For certain linear scattering or radiation problems, a known background Green's function represents the field through equivalent boundary sources or currents. The unknown changes from a volume field to a boundary density determined by boundary conditions. This can reduce dimension, but the long-range Green's function generally produces a dense matrix.

\textbf{Core Conceptual Connection:}
\[
\text{field PDE}
\longrightarrow
\text{integral equation}
\longrightarrow
\text{current expansion}
\]
\end{authornote}

and constrain the integral-equation residual with test functions \(w_m\):
\[
\avg{w_m,\Lop J}
=
\avg{w_m,g}.
\]

\begin{askai}
Ask AI:
\begin{quote}
Discretize a simple one-dimensional integral equation using three piecewise-constant basis functions:
\[
\Lop J=g
\]
Write the entries \(\avg{w_m,\Lop f_n}\) explicitly and explain why each test function generally couples to every basis function.
\end{quote}

Follow-up:
\begin{quote}
Compare the nonzero patterns of differential- and integral-equation matrices, and identify the costs targeted by fast multipole methods or hierarchical matrices.
\end{quote}
\end{askai}

This gives
\[
\bm Z\bm I=\bm V.
\]

The long-range kernel generally makes $\bm Z$ dense. This structure and fast matrix application lead to Volume~0.

Specifically, if $(\mathcal LJ)(x)=\int_\Gamma G(x,y)J(y)\dd S_y$,
then
\[
Z_{mn}=\int_\Gamma w_m(x)\int_\Gamma G(x,y)f_n(y)\dd S_y\dd S_x,\qquad
V_m=\int_\Gamma w_mg\dd S.
\]
This testing pairing belongs to integral-equation discretization. Stability depends on the operator, trace spaces, kernel, and basis.

\part{Modern High-Order Discretization}
\chapter{Discontinuous Galerkin Methods}\label{ch:10}
\begin{chaptersummary}
Finite volumes enforce intercell conservation through shared fluxes; finite elements constrain residuals through testing. DG combines both mechanisms with local high-order polynomials and elementwise weak forms. One-dimensional advection gives the mass matrix and local evolution, while tests with constants and the solution yield conservation and energy identities. Elliptic penalties then address diffusion interfaces.
\end{chaptersummary}
\learninggoals{Construct broken trial spaces and traces; derive one-dimensional DG matrices; prove local and global conservation; distinguish hyperbolic upwind stability from elliptic penalty coercivity.}
\studyroute{General nonlinear entropy stability, multidimensional boundary estimates, and elliptic DG error theorems are omitted. The derivation assumes constant coefficients, affine elements, and exact integration.}{discontinuous Galerkin; broken polynomial space; upwind flux; symmetric interior penalty; modal nodal DG}{For a quadratic Legendre basis, assemble the advection mass matrix, volume matrix, and endpoint vectors, then sum interface contributions to verify the energy estimate.}

\section{Local Spaces, Degrees of Freedom, and Traces}\index{DG methods}
On a periodic interval, let $K_j=(x_{j-1/2},x_{j+1/2})$ have length $h_j$.
Take
\[
V_h^p=\{v\in L^2:v|_{K_j}\in\Poly_p(K_j)\}.
\]
Each element has $p+1$ scalar degrees of freedom; neighboring traces are independent. Consider
\[
u_t+a u_x=0,\qquad F(u)=au,\qquad a\in\R.
\]
Fix the interface orientation from left to right, with traces $u^-$ and $u^+$,
$[u]=u^+-u^-$, and $\{u\}=(u^-+u^+)/2$. A consistent upwind flux is
\[
\widehat f=a\{u\}-\frac{|a|}{2}[u].
\]
For $a>0$, $\widehat f=au^-$.

\conceptbridge{When elementwise approximations may jump, how does information cross interfaces?}
{The broken space permits independent element approximations. Integration by parts generates interface fluxes; a shared Riemann flux couples the two traces and their local residuals. Constant test functions give cell conservation, while higher-order test functions determine the evolution of the internal polynomial.}
{Broken polynomials $\longrightarrow$ elementwise weak form $\longrightarrow$ shared numerical flux.}

\section{The Complete Weak Form for One-Dimensional Advection}
For any $v_h\in\Poly_p(K_j)$, multiply the equation and integrate by parts:
\[
\int_{K_j}v_h\partial_tu_h\dd x
-\int_{K_j}a u_h\partial_xv_h\dd x
+\widehat f_{j+1/2}v_h(x_{j+1/2}^{-})
-\widehat f_{j-1/2}v_h(x_{j-1/2}^{+})=0.
\]
Map $x=x_j+J_j\xi$, with $\xi\in[-1,1]$ and $J_j=h_j/2$.
Expand $u_h=\sum_{\ell=0}^pU_{j\ell}\phi_\ell(\xi)$ and define
\[
M_{i\ell}=\int_{-1}^1\phi_i\phi_\ell\dd\xi,\quad
S_{i\ell}=\int_{-1}^1\phi_i'\phi_\ell\dd\xi,\quad
r_i=\phi_i(1),\quad l_i=\phi_i(-1).
\]
Testing with each basis function gives
\begin{equation}\label{eq:dg-local}
J_jM\dot U_j=a S U_j-r\widehat f_{j+1/2}+l\widehat f_{j-1/2}
=:R_j(U).
\end{equation}
The volume matrix acts only within the element; interface residuals use neighboring traces.
In a Legendre modal basis,
$M_{ii}=2/(2i+1)$. For $p=1$,
\[
M=\begin{pmatrix}2&0\\0&2/3\end{pmatrix},\quad
S=\begin{pmatrix}0&0\\2&0\end{pmatrix},\quad
r=\begin{pmatrix}1\\1\end{pmatrix},\quad
l=\begin{pmatrix}1\\-1\end{pmatrix}.
\]
With $a>0$ and $Q_j=U_{j0}+U_{j1}$,
\[
h_j\dot U_{j0}=-aQ_j+aQ_{j-1},\qquad
\frac{h_j}{3}\dot U_{j1}=2aU_{j0}-aQ_j-aQ_{j-1}.
\]
This is a directly implementable local ODE for the degrees of freedom.

\section{Local Conservation and the Energy Identity}
Taking $v_h=1$ gives
\[
\frac{\dd}{\dd t}\int_{K_j}u_h\dd x
=-\widehat f_{j+1/2}+\widehat f_{j-1/2}.
\]
Shared fluxes cancel when summed over all cells. Upwind DG with $p=0$ recovers first-order finite-volume spatial discretization.
Taking $v_h=u_h$ and combining the two sides of each interface gives
\[
\frac12\frac{\dd}{\dd t}\sum_j\int_{K_j}u_h^2\dd x
=\sum_f\left(\widehat f-a\{u_h\}\right)[u_h]
=-\frac{|a|}{2}\sum_f[u_h]^2.
\]
Periodicity, constant coefficients, and exact integration close this identity. A central flux gives a zero right-hand side.
Fully discrete stability additionally requires compatibility between the time integrator and spatial spectrum. Physical boundaries contribute boundary fluxes to the energy identity.

\section{Modal and Nodal Forms; Elliptic Problems}
Modal DG stores basis coefficients; nodal DG stores values at interpolation nodes.
If both represent the same polynomial space with identical integration and fluxes, a basis transformation connects their semidiscrete equations.
Mass lumping or altered quadrature can change the discrete operator.

For periodic one-dimensional Poisson, the symmetric interior penalty form can be used on the mean-zero broken space:
\begin{equation}\label{eq:sip}
\begin{aligned}
b_h(u,v)={}&\sum_K\int_Ku'v'\dd x
+\sum_f\bigl(\{u'\}[v]+\{v'\}[u]\bigr)\\
&+\sum_f\frac{\sigma(p+1)^2}{h_f}[u][v].
\end{aligned}
\end{equation}
The scalar jump remains right trace minus left trace, so the consistency terms have positive signs.
Choose $h_f$ as an appropriate minimum of neighboring sizes. For sufficiently large $\sigma$,
inverse trace estimates absorb the interface cross terms into the volume and penalty terms, giving discrete coercivity.
Dirichlet boundaries require compatible boundary consistency and penalty terms; see~\cite{arnold2002}.

\aiexercise{For a uniform periodic mesh with $p=2$, calculate $M,S,l,r$ from~\eqref{eq:dg-local} and construct the global block coupling. Test with constants and the solution itself to verify total mass conservation and the upwind energy identity.}
{Insert a smooth exact solution into~\eqref{eq:sip} and check all interface signs. Use a discrete trace inequality to identify the cross terms controlled by the penalty parameter.}

\chapter{Flux Reconstruction}\label{ch:11}
\begin{chaptersummary}
The strong DG form represents interface residuals as liftings into the element interior. FR follows this idea by interpolating the local flux and incorporating a shared flux through endpoint correction functions. This chapter derives the nodal update and conservation property, then identifies the exact relationship with a particular DG discretization through a mass-matrix lifting condition.
\end{chaptersummary}
\learninggoals{Construct discontinuous and corrected fluxes for a one-dimensional conservation law; write the FR spatial operator; prove cell-average conservation; state the basis, integration, and flux conditions required for a DG equivalence.}
\studyroute{Multidimensional curved meshes, nonlinear aliasing, and general FR entropy analysis are omitted. The next chapter constructs energy-stable correction functions.}{Huynh flux reconstruction; correction function; lifting operator; SBP correction procedure}{At three Gauss nodes, construct a quadratic solution and cubic correction functions, assemble the nodal FR operator, and compare it with exactly integrated strong-form DG.}

\section{Discontinuous Fluxes and Endpoint Corrections}\index{FR methods}
Consider $u_t+\partial_xF(u)=0$ on an affine element $x=x_j+J_j\xi$.
Choose $p+1$ distinct nodes $\xi_i$ with Lagrange basis $\ell_i$:
\[
u_j^\delta(\xi)=\sum_{i=0}^pu_{ji}\ell_i(\xi),\qquad
f_j^D(\xi)=\sum_{i=0}^pF(u_{ji})\ell_i(\xi).
\]
Here $f^D$ is the degree-$p$ flux interpolant; the nonlinear flux $F(u^\delta)$ may have higher degree.
An interface solver uses the two solution traces to produce shared, consistently oriented fluxes $f_L^*,f_R^*$.
Choose correction functions of degree at most $p+1$ satisfying
\[
g_L(-1)=1,\quad g_L(1)=0,\qquad
g_R(-1)=0,\quad g_R(1)=1.
\]
The corrected flux is
\begin{equation}\label{eq:fr-flux}
f^C=f^D+(f_L^*-f^D(-1))g_L+(f_R^*-f^D(1))g_R.
\end{equation}
Thus $f^C(-1)=f_L^*$ and $f^C(1)=f_R^*$.
The globally shared interface flux has the same endpoint value in each local polynomial. For the original FR construction, see~\cite{huynh2007}.

\conceptbridge{How can two endpoint flux values change every interior degree of freedom?}
{Correction functions extend endpoint flux differences into element-interior polynomials. Their derivatives deliver interface residuals to every solution node. The correction shape determines how this lifting acts on polynomial modes and therefore affects the spectrum and stability.}
{Endpoint flux differences $\longrightarrow$ correction polynomials $\longrightarrow$ interior residual lifting.}

\section{The Nodal FR Update and Local Conservation}
Let $D_{ik}=\ell_k'(\xi_i)$, $C_{iL}=g_L'(\xi_i)$, and $C_{iR}=g_R'(\xi_i)$.
Constrain the corrected differential residual at the nodes:
\begin{equation}\label{eq:fr-operator}
J_j\dot u_{ji}
=-\sum_kD_{ik}F(u_{jk})
-(f_L^*-f_L^D)g_L'(\xi_i)
-(f_R^*-f_R^D)g_R'(\xi_i).
\end{equation}
Let $R$ be the $2\times(p+1)$ endpoint-evaluation matrix. Then
$J_j\dot{\bm u}_j=-D\bm f_j-C(\bm f^*-R\bm f_j)$.
All derivatives have degree at most $p$, so the nodal equality is also a polynomial identity throughout the element.
Exact integration gives
\[
\frac{\dd}{\dd t}\left(J_j\int_{-1}^1u_j^\delta\dd\xi\right)
=-f_R^*+f_L^*.
\]
Shared fluxes and endpoint conditions therefore ensure cell-average conservation. The particular correction functions must additionally satisfy the desired stability conditions.

\section{The Matrix Relationship with DG}
In the same basis, let $M_{ik}=\int\ell_i\ell_k$ and $B=\diag(-1,1)$.
Polynomial integration by parts gives
\[
MD+D^{\mathsf T}M=R^{\mathsf T}BR.
\]
For a linear flux $F(u)=au$, strong-form DG is
\[
J_j\dot{\bm u}_j=-aD\bm u_j-M^{-1}R^{\mathsf T}B
(\bm f^*-aR\bm u_j).
\]
Therefore, when
\begin{equation}\label{eq:fr-dg-lift}
C=M^{-1}R^{\mathsf T}B,
\end{equation}
and representation, integration, mapping, and fluxes agree, FR and this semidiscrete DG operator coincide.
An invertible basis transformation gives the same relationship in a modal implementation.
Nonlinear fluxes, underintegration, or other correction functions require a new comparison of volume and boundary operators.

\section{Error and Spectral Analysis}
Correction functions change the interface feedback into element modes. On a uniform periodic mesh, the Bloch symbol gives the spectrum,
allowing comparison of dissipation, dispersion, and explicit stable steps on the physical branch.
An energy method instead examines evolution directly in a chosen norm and includes physical boundaries.
Dissipation, dispersion, and stable step size also depend on fluxes, correction functions, and the time integrator. Cost comparisons by total degrees of freedom and error therefore record these choices alongside formal order.

\aiexercise{For $p=1$, a linear flux, and Gauss solution nodes, calculate $M,D,R,B$. Obtain $g_L',g_R'$ from~\eqref{eq:fr-dg-lift}, integrate them, check all four endpoint conditions, and compare the local matrices with DG.}
{Replace a correction function by another polynomial satisfying the same endpoint conditions. Check conservation, then calculate the periodic-symbol eigenvalues to distinguish the requirements for conservation and stability.}

\chapter{VCJH Correction Functions and Energy Stability}\label{ch:12}
\begin{chaptersummary}
FR endpoint conditions ensure conservation but leave freedom in the interior correction. A modified energy norm constrains this freedom for one-dimensional linear advection, producing the VCJH one-parameter family. With normalization, mass matrix, and quadrature fixed, its relationships with DG, particular SD constructions, and Huynh corrections can be stated precisely.
\end{chaptersummary}
\learninggoals{Write the parameterized VCJH corrections; derive positivity and the linear energy identity; identify DG and SD parameters under specified assumptions; assess parameters through spectra and errors.}
\studyroute{The proof assumes one-dimensional constant-coefficient advection, stationary affine meshes, and exact polynomial integration. Nonlinear entropy, curved meshes, and diffusion extensions are further reading.}{VCJH correction function; energy stable flux reconstruction; modified Sobolev norm; SBP CPR}{For $p=1,2,3$, calculate correction endpoints, lifting matrices, and modified mass matrices, and verify the linear energy matrix identity.}

\section{Parameterized Correction Functions}\index{VCJH methods}
Let $p\ge1$, with standard Legendre polynomial $P_p$,
whose leading coefficient is $a_p=(2p)!/[2^p(p!)^2]$.
Use the normalization of~\cite{vincent2011}:
\[
\eta_p=\frac{c(2p+1)(a_pp!)^2}{2},
\]
\begin{equation}\label{eq:vcjh-functions}
\begin{aligned}
g_L(\xi)&=\frac{(-1)^p}{2}
\left[P_p(\xi)-\frac{\eta_pP_{p-1}(\xi)+P_{p+1}(\xi)}{1+\eta_p}\right],\\
g_R(\xi)&=\frac12
\left[P_p(\xi)+\frac{\eta_pP_{p-1}(\xi)+P_{p+1}(\xi)}{1+\eta_p}\right].
\end{aligned}
\end{equation}
The identities $P_k(1)=1$ and $P_k(-1)=(-1)^k$ give the four FR endpoint conditions directly.
The parameter $c$ is dimensionless. Changing it preserves flux consistency while modifying the lifting into interior modes.

\conceptbridge{Why introduce a modified energy with higher derivatives rather than check endpoint conditions alone?}
{Endpoint conditions constrain only the integrated flux. An energy estimate must also control the pairing of interior modes with interface liftings. Adjusting the norm of the highest mode permits a lifting satisfying discrete integration by parts; positivity ensures that the estimate controls every degree of freedom.}
{Conservative endpoint conditions $\longrightarrow$ interior lifting $\longrightarrow$ a positive modified energy.}

\section{Modified Norm and Stable Parameter Range}
In an exactly integrated polynomial basis, let $D$ be the reference derivative matrix and $M$ the mass matrix. Define
\begin{equation}\label{eq:vcjh-norm}
K_c=\frac c2(D^p)^{\mathsf T}MD^p,\qquad H_c=M+K_c.
\end{equation}
Then
\[
U^{\mathsf T}H_cU
=\int_{-1}^1u_h^2\dd\xi
+\frac c2\int_{-1}^1(\partial_\xi^pu_h)^2\dd\xi.
\]
In a Legendre modal basis, only the highest-order diagonal entry changes:
\[
H_c=\diag\left(2,\frac23,\ldots,
\frac{2}{2p+1}+c(a_pp!)^2\right).
\]
Hence
\begin{equation}\label{eq:vcjh-bound}
c>c_-:=-\frac{2}{(2p+1)(a_pp!)^2}
\quad\Longleftrightarrow\quad H_c>0
\quad\Longleftrightarrow\quad\eta_p>-1.
\end{equation}
Since $D^{p+1}=0$, $K_cD=0$. The derivative matrix of the corrections~\eqref{eq:vcjh-functions} satisfies
\[
C=H_c^{-1}R^{\mathsf T}B,\qquad
H_cD+D^{\mathsf T}H_c=R^{\mathsf T}BR.
\]
This matrix interpretation is compatible with SBP--CPR constructions~\cite{ranocha2016}.

\section{The Linear Energy Proof}
For constant $a$, multiply the FR equation by $U_j^{\mathsf T}H_c$.
Discrete integration by parts converts the volume term to endpoint trace squares, and the lifting inserts the shared flux.
Summation over periodic elements gives
\[
\frac{\dd}{\dd t}\sum_j\frac{J_j}{2}U_j^{\mathsf T}H_cU_j
=\sum_f(\widehat f-a\{u_h\})[u_h]
=-\frac{|a|}{2}\sum_f[u_h]^2.
\]
The constant mode satisfies $K_c\bm1=0$, so this lifting also preserves the physical cell average.
Negative parameters within~\eqref{eq:vcjh-bound} still define a positive norm.
Near the lower bound, the highest-mode weight decreases; norm-equivalence constants, spectral radius, and errors must therefore also be assessed.

\section{Special Parameters and Practical Selection}
For the exact mass matrix and linear-advection setting used here:
\begin{center}\small
\begin{tabularx}{\textwidth}{@{}l>{\raggedright\arraybackslash}X@{}}
\toprule
Parameter & Specific relationship\\
\midrule
$c=0$ & $H_c=M$, recovering the exactly integrated DG lifting of Chapter~\ref{ch:11}.\\
$\eta_p=p/(p+1)$ & A particular one-dimensional SD construction with interior flux points at the roots of $P_p$ and interface endpoints.\\
$\eta_p=(p+1)/p$ & Huynh's $g_2$ correction; its relationship with a particular Lobatto-quadrature DGSEM additionally requires fixed quadrature and nodes.\\
\bottomrule
\end{tabularx}
\end{center}
These relationships use the specified correction functions, nodes, and quadrature.
Changes in SD nodes, nonlinear flux interpolation, or dimension can be assessed through the same volume and lifting matrices.
If a reference uses $\kappa$ with $K=\kappa(D^p)^{\mathsf T}MD^p$,
then the present parameter is $c=2\kappa$.
Parameter selection can compare dispersion at target wave numbers, available CFL numbers, and error tolerances. The linear energy range provides an analytically controlled starting point.

\aiexercise{In a $p=2$ Legendre basis, construct $D,M,R,B$. Calculate $H_c,C$ for $c=0$, the SD parameter, and a negative stable parameter. Verify positivity, $H_cC=R^{\mathsf T}B$, and the modified-energy derivative.}
{At fixed total degrees of freedom, upwind flux, and RK4 integration, measure smooth-wave errors and stable steps for these parameters. Report semidiscrete and fully discrete results separately, and identify nonlinear stability as further reading.}

\chapter{Mortar Methods and Nonconforming Mesh Coupling}\label{ch:13}
\begin{chaptersummary}
Local high-order spaces allow independent element sizes and degrees, so interface nodes and trace spaces may not match. Mortars organize projection, flux evaluation, and return transfer on a shared interface, ensuring that both sides use the same integral data. Projection matrices, constant preservation, and linear energy conditions provide an interface for hp adaptation and heterogeneous methods.
\end{chaptersummary}
\learninggoals{Construct a common interface partition and mortar space; calculate trace projection and dual return transfer; prove cancellation of integrated fluxes; state the quadrature and trace-reproduction conditions for energy stability.}
\studyroute{Inconsistent curved geometry, nonlinear entropy projections, and specialized reduced-order mortar stabilization are omitted. The linear proof uses a common space reproducing both traces.}{mortar method; nonconforming interface; mass-adjoint projection; hp nonconforming DG}{On one segment, construct mortar mass and projection matrices for a linear left trace and constant right trace, then check flux loads and constant tests explicitly.}

\section{Common Interfaces and Trace Spaces}\index{Mortar methods}
Consider a planar edge $\Gamma$ between two two-dimensional elements, with normal fixed from left to right.
For unequal polynomial degrees, or one coarse edge adjoining several fine edges, first partition the physical interface into common subfaces.
On each subface, choose $V_m=\Poly_{p_m}$ and use identical geometry mappings, normals, and integration measures.
Taking $p_m$ at least as large as both trace degrees reproduces both traces on affine subfaces.
For this conservative coupling approach, see~\cite{kopriva1996}.

\conceptbridge{Why introduce a common interface instead of evaluating fluxes at each side's own nodes?}
{Different nodes lead to different flux samples and integrals. A common interface represents both traces in one space,
evaluates one numerical flux, and returns it with opposite outward normals. Constant preservation ensures agreement of total quantities;
compatible integration and dual transfer also connect the two energy pairings.}
{Nonmatching traces $\longrightarrow$ common interface space $\longrightarrow$ shared flux and compatible return transfer.}

\section{Projection, Fluxes, and Dual Return Transfer}
Let $\mu_\alpha$ be a mortar basis and $\phi_i^s$ a trace basis on side $s=L,R$. Define
\[
(M_m)_{\alpha\beta}=\int_\Gamma\mu_\alpha\mu_\beta\dd S,\qquad
(B_{ms})_{\alpha i}=\int_\Gamma\mu_\alpha\phi_i^s\dd S.
\]
The $L^2$ trace projection satisfies
\begin{equation}\label{eq:mortar-proj}
P_s=M_m^{-1}B_{ms},\qquad u_{s,m}=P_su_s.
\end{equation}
Evaluate a single-valued normal flux from $u_{L,m},u_{R,m}$ at common quadrature nodes.
Represent its load with the same rule as $M_mf_m^*$. The outward flux loads are
\begin{equation}\label{eq:mortar-return}
r_L=P_L^{\mathsf T}M_mf_m^*,\qquad
r_R=-P_R^{\mathsf T}M_mf_m^*.
\end{equation}
These are dual loads in weak residuals. Flux coefficients on one side require an additional inverse trace mass matrix;
this also gives the mass-adjoint relation $M_sP_{m\to s}=P_s^{\mathsf T}M_m$.

\section{Conservation and Linear Energy Conditions}
Let $\bm1_s$ and $\bm1_m$ represent constants. Constant reproduction means
$P_s\bm1_s=\bm1_m$, hence
\[
\bm1_L^{\mathsf T}r_L+\bm1_R^{\mathsf T}r_R=0.
\]
If one coarse edge meets several mortar subfaces, sum their loads;
the coarse and fine sides still use the same integrals.

For a linear energy estimate, further assume that the common space reproduces both traces exactly
and that quadrature exactly integrates the relevant products. Then
$P_s^{\mathsf T}M_mP_s=M_s$,
For multiple subfaces, sum the matrices on the left.
For constant normal speed $a_n$, use
\[
f_m^*=a_n\{u_m\}-\frac{|a_n|}{2}[u_m],
\]
The combined DG interface energy contribution is
\[
-\frac{|a_n|}{2}[u_m]^{\mathsf T}M_m[u_m]\le0.
\]
If the common space has lower degree or the geometry and quadrature differ, the trace-norm identity requires a corresponding replacement construction.
Constant preservation supplies total conservation; trace reproduction, quadrature, and mass-adjoint transfer additionally support the energy estimate.

\section{A Hand-Computable Nonconforming Interface}
Let $\Gamma=[-1,1]$, with left trace basis $(1,s)$ and a constant-only right trace.
Use mortar basis $(1,s)$. Exact integration gives
\[
M_m=\begin{pmatrix}2&0\\0&2/3\end{pmatrix},\quad
P_L=I_2,\quad P_R=\begin{pmatrix}1\\0\end{pmatrix}.
\]
For shared flux coefficients $(F_0,F_1)^{\mathsf T}$,
\[
r_L=\begin{pmatrix}2F_0\\2F_1/3\end{pmatrix},\qquad r_R=-2F_0.
\]
a constant test reads only the first left component, so the total flux cancels exactly.
The linear component contributes to the left trace residual. Changing interface spaces during hp adaptation requires updating projections and physical measures together.

\aiexercise{Use a quadratic left trace and linear right trace, with three Gauss points for exact integration. Construct $M_m,B_{mL},B_{mR}$ and verify constant reproduction, mass adjointness, and upwind interface energy.}
{Split the right edge into two subfaces with linear traces. Write the subface-map Jacobians and prove that the coarse-side total load is the negative sum of the two fine-side loads.}

\chapter{Combining High-Order Methods and Time Integration}\label{ch:14}
\begin{chaptersummary}
DG, FR, and mortars define local spaces and interface operators; complete evolution also requires temporal coordination. Shared fluxes connect spatial methods, ADER and CERK provide within-step states, LTS coordinates local steps, and IMEX treats stiff splitting. Spatial approximation, interface integration, and time advancement therefore form one accuracy and conservation problem.
\end{chaptersummary}
\learninggoals{Construct shared-flux DG--FR coupling; distinguish ADER, CERK, and LTS; write an IMEX split; design conservative h, p, and hp transfers and time-convergence experiments.}
\studyroute{High-order RK coefficient derivations, general nonlinear space-time predictor iteration, and adaptive-error proofs are omitted. Each algorithm must specify its quadrature and time schedule.}{ADER DG; continuous explicit Runge--Kutta; conservative local time stepping; IMEX; hp adaptivity}{Couple neighboring DG and FR elements with one interface solver. Verify conservation at a common step size, then introduce two-to-one LTS and compare interface time integrals and total conservation errors.}

\section{Shared Fluxes between Different Spatial Methods}
For a common conservation law, both DG and FR express cell-average updates through outward flux integrals.
Using common states, normals, and single-valued fluxes at the interface, returned through compatible liftings,
allows a DG element to couple to an FR element. Nonconforming interfaces use the mortars of Chapter~\ref{ch:13}.
Overall order is constrained by both approximations, interface projection, and quadrature;
the individual energy estimates must also close under the common interface pairing.

\conceptbridge{If each spatial method is stable, why recheck interfaces and time integration after coupling?}
{Stability estimates contain interface and time-update contributions. Coupling changes the times, projections, and pairings of neighboring data,
and therefore changes these contributions. Shared flux integrals control conservation; stability regions and compatible energy pairings control evolution.
Both checks must apply to the final combined operator.}
{Local spatial operators $\longrightarrow$ interface and temporal coupling $\longrightarrow$ fully discrete structure.}
\section{High-Order Time Integration: ADER, CERK, and LTS}\label{sec:v1-time}
Spatial discretization gives the finite-dimensional system
\[
M_h\dot U_h=R_h(U_h,t).
\]
For a fixed invertible mass matrix, write $\dot U_h=L_h(U_h,t)$.
Time integrators approximate its evolution; their order, stability region, and interface time data must match the spatial discretization.
The following distinguishes within-step approximations from local time scheduling. Volume~2, Section~5.5 develops the Maxwell setting.

\paragraph{ADER: Local Space-Time Prediction and Conservative Updating.}
ADER (Arbitrary high-order DERivatives) originated in high-order time expansions:
classical constructions replace temporal derivatives with spatial derivatives through the PDE, while modern constructions can solve an element-local space-time weak problem.
The resulting $q_K(\bm x,t)$ supplies volume terms and interface flux integrals throughout the time step.
For a source-free finite-volume update,
\[
|K|(\bar U_K^{n+1}-\bar U_K^n)
=-\sum_{f\subset\partial K}\int_{t_n}^{t_{n+1}}\int_f
\widehat F(q_K,q_{K_f};\bm n_f)\dd S\dd t.
\]
the local prediction is followed by a correction coupling neighbors through shared fluxes. Accuracy depends on reconstruction, the predictor, and integration order;
explicit hyperbolic updates retain their corresponding CFL constraint.
See \href{https://doi.org/10.1016/j.cma.2014.07.019}{Dumbser's ADER--LTS construction}.

In ADER--DG, the same local predictor $q_K$ also supplies basis-weighted volume residuals.
Integrating the DG weak form over one step gives
\[
\begin{aligned}
\int_Kv_h(u_h^{n+1}-u_h^n)\dd x
={}&\int_{t_n}^{t_{n+1}}\int_K\nabla v_h\cdot F(q_K)\dd x\dd t\\
&-\int_{t_n}^{t_{n+1}}\int_{\partial K}v_h\widehat F(q_K,q_{K_f};n)\dd S\dd t.
\end{aligned}
\]
The predictor supplies space-time data; DG testing and shared fluxes form the correction. This identifies precisely where ADER couples to spatial DG.

\paragraph{CERK: Continuous Explicit Runge--Kutta.}
CERK constructs a continuous within-step approximation from explicit RK stage derivatives $k_i$:
\[
U_h(t_n+\theta\Delta t)\approx
U_h^n+\Delta t\sum_{i=1}^{s}b_i(\theta)k_i,\qquad 0\le\theta\le1.
\]
Here $b_i(0)=0$ and $b_i(1)=b_i$, matching the endpoint RK weights.
Continuous-extension order conditions control within-step accuracy, while the endpoint RK conditions control end-of-step accuracy. Coupling and output must satisfy their respective requirements.
The approximation supports output, coupling interpolation, or local predictors. See
\href{https://doi.org/10.1137/0913084}{Owren--Zennaro's continuous explicit RK methods}.
For element-local prediction, RK stages solve a local problem and neighboring effects enter through correction fluxes;
see~\cite{altmann2010} for this space-time coupling.

\paragraph{LTS: Local Time Stepping.}
When local mesh sizes or wave speeds vary strongly, the most restrictive element determines a uniform time step.
LTS assigns different physical time steps to different elements and synchronizes them at common time levels.
If a coarse element advances by $\Delta T$ while a fine neighbor takes $m$ substeps, their shared time-integrated flux must satisfy
\[
\Phi_f^{\mathrm{coarse}}
=\sum_{r=1}^{m}\Phi_{f,r}^{\mathrm{fine}},
\]
All terms use one normal convention and enter neighboring updates with opposite signs.
ADER or CERK can provide within-step states; LTS additionally requires scheduling, quadrature, and flux accumulation.
Physical-time LTS must preserve temporal accuracy; local pseudo-time stepping for steady-state acceleration has a different objective.
Efficiency depends on the fraction of small elements and the costs of interface prediction, integration, and synchronization.

\begin{askai}
\textbf{Ask AI:}\par
\textbf{Omitted Detail and Keywords:} ADER predictor--corrector; CERK dense-output order conditions; conservative LTS. High-order coefficient construction and scheduling algorithms are omitted.

For one-dimensional advection, place coarse and fine elements next to each other with time-step ratio $2:1$. Use a time polynomial and a CERK extension to supply interface states, list the flux integrals over both fine substeps, and verify total conservation.

Follow-up: Freeze the coarse-side state for the whole step and measure the resulting temporal order; then compare the update counts of uniform stepping and LTS.
\end{askai}

\section{RK Stability Regions and IMEX Splitting}
If a spatial mode satisfies $\dot y=\lambda_hy$, an RK update is
$y^{n+1}=\mathcal R(\Delta t\lambda_h)y^n$,
The spatial spectrum and temporal stability region must be checked together.
SSP methods express updates as convex combinations of suitable forward Euler steps.
Inheritance of strong stability requires the Euler step to be stable in the target functional under the corresponding step restriction.
High-order spatial energy estimates and SSP properties have distinct assumptions.

For combined propagation and stiff relaxation, use the split
\[
\dot U=F_h(U)+G_h(U)
\]
with explicit $F_h$ and implicit $G_h$.
The common stages of an autonomous additive Runge--Kutta method are
\[
\begin{aligned}
Y_i&=U^n+\Delta t\sum_{j<i}\widetilde a_{ij}F_h(Y_j)
+\Delta t\sum_{j\le i}a_{ij}G_h(Y_j),\\
U^{n+1}&=U^n+\Delta t\sum_i\widetilde b_iF_h(Y_i)
+\Delta t\sum_i b_iG_h(Y_i).
\end{aligned}
\]
The two tableaux must satisfy additive order conditions~\cite{ascher1997}.
Implicit treatment of local collision sources can reduce the step restriction imposed by fast relaxation.
Explicit propagation retains its stability restriction, and the stiff limit requires checks of asymptotic preservation and actual temporal order.
Volume~2 gives a conducting-Maxwell example.

\section{h, p, and hp Adaptation}
$h$ adaptation changes element sizes; $p$ adaptation changes local degrees; $hp$ adaptation changes both.
Smooth regions may benefit from degree elevation, whereas discontinuities and local geometric scales often require smaller elements and stabilization.
Transfer operators should satisfy
\[
\int_Ku_{\mathrm{new}}\dd x
=\int_Ku_{\mathrm{old}}\dd x,
\]
Sum over all children during refinement; aggregate child integrals before coarsening.
An $L^2$ projection preserves the average when constants are included as tests; positivity and other constraints may require additional limiting.
After repartitioning, update interface spaces, geometric measures, and quadrature rules together.

\aiexercise{For smooth periodic advection, fix the spatial mesh and measure temporal convergence with RK4 and a Heun continuous extension. Then fix a sufficiently small step and measure spatial convergence. Explain what mass, energy, and phase errors each measure.}
{For a discrete-velocity equation with fast BGK relaxation, select documented IMEX coefficients and distinguish propagation CFL, collision relaxation, and an unresolved limiting regime. Record the coefficient source and observed order as part of the experiment.}

\part{Mesoscopic Kinetics and LBM}
\chapter{From the Boltzmann Equation to Lattice Boltzmann Methods}\label{ch:15}
\begin{chaptersummary}
The preceding methods discretize macroscopic equations directly. This chapter instead selects discrete velocities and collisions at the kinetic level, recovering macroscopic transport in a long-wave, low-Mach limit. D2Q9 and D3Q27 connect velocity sets, isotropic moments, and exact lattice streaming; Chapman--Enskog expansion links nonequilibrium stress and the discrete-time correction to viscosity.
\end{chaptersummary}
\learninggoals{Define discrete populations and conserved moments; construct D2Q9 and D3Q27; perform streaming and collision updates; explain the low-Mach fluid limit and viscosity formula through a multiscale expansion.}
\studyroute{Full third-order Chapman--Enskog algebra and general thermal-lattice construction are omitted. The model is standard, isothermal, low Mach, and unforced.}{lattice Boltzmann; D2Q9; D3Q27; Chapman--Enskog expansion; lattice viscosity}{Verify second- and fourth-order isotropy of D2Q9, then derive the half-time-step viscosity correction from the second-order Taylor expansion of streaming and collision.}

\section{Discrete Velocities and Mesoscopic Variables}\index{LBM}
Continuous kinetics contains both coordinate sets $\bm x,\bm v$.
A discrete-velocity model selects velocities $\bm c_i$ and populations $f_i(\bm x,t)$ satisfying
\[
\partial_tf_i+\bm c_i\cdot\nabla f_i=\Omega_i,\qquad
\rho=\sum_i f_i,\qquad \rho\bm u=\sum_i\bm c_i f_i.
\]
Here $f_i$ are mass-density populations with velocity-quadrature weights absorbed,
unlike the continuous number-density distribution $f$ in Chapter~1.
Moment matching by the velocity set and equilibrium determines the macroscopic models that can be represented;
spatial streaming and time discretization then define the numerical update.
If $\bm c_i\Delta t$ connects lattice nodes exactly, characteristic streaming becomes a direct shift of lattice data.

\conceptbridge{How can finitely many velocity directions recover macroscopic viscous flow?}
{Macroscopic equations read low-order distribution moments. If velocities, weights, and equilibrium reproduce the required isotropic moments,
and collisions conserve mass and momentum, a long-wave expansion determines the nonequilibrium stress.
The recovered equations and their errors depend jointly on these moment conditions, scale ordering, and time discretization.}
{Discrete velocities and collision moments $\longrightarrow$ nonequilibrium stress $\longrightarrow$ macroscopic transport limit.}

\section{D2Q9 and D3Q27}
Let $c_\ell=\Delta x/\Delta t$. D2Q9 uses velocities
\[
\bm c_i/c_\ell\in\{(0,0),(\pm1,0),(0,\pm1),(\pm1,\pm1)\}.
\]
The rest, axial, and diagonal weights are
$4/9,1/9,1/36$, respectively, with lattice sound speed $c_s^2=c_\ell^2/3$.
D3Q27 is the tensor product of $\{0,\pm c_\ell\}$ in three directions,
with one-dimensional weights $w(0)=2/3$, $w(\pm c_\ell)=1/6$; three-dimensional weights are their products.
Both lattices satisfy, in their respective dimensions,
\[
\sum_iw_i=1,\quad \sum_iw_i c_{i\alpha}=0,\quad
\sum_iw_i c_{i\alpha}c_{i\beta}=c_s^2\delta_{\alpha\beta},
\]
\[
\sum_iw_i c_{i\alpha}c_{i\beta}c_{i\gamma}c_{i\delta}
=c_s^4(\delta_{\alpha\beta}\delta_{\gamma\delta}
+\delta_{\alpha\gamma}\delta_{\beta\delta}
+\delta_{\alpha\delta}\delta_{\beta\gamma}).
\]
Higher-order velocity moments are constrained by the finite lattice.

\section{Equilibrium, Collision, and Streaming}
The second-order isothermal low-Mach equilibrium is
\begin{equation}\label{eq:lbm-eq}
f_i^{\mathrm{eq}}=w_i\rho\left[
1+\frac{\bm c_i\cdot\bm u}{c_s^2}
+\frac{(\bm c_i\cdot\bm u)^2}{2c_s^4}
-\frac{|\bm u|^2}{2c_s^2}\right].
\end{equation}
The moment identities give
\[
\sum_if_i^{\mathrm{eq}}=\rho,\quad
\sum_i\bm c_if_i^{\mathrm{eq}}=\rho\bm u,\quad
\sum_i\bm c_i\bm c_if_i^{\mathrm{eq}}=\rho c_s^2I+\rho\bm u\bm u.
\]
Each BGK--LBM step performs local collision followed by streaming:
\begin{equation}\label{eq:lbe}
f_i^*(\bm x,t)=f_i-\frac{\Delta t}{\tau_{\mathrm{LB}}}(f_i-f_i^{\mathrm{eq}}),
\qquad
f_i(\bm x+\bm c_i\Delta t,t+\Delta t)=f_i^*(\bm x,t).
\end{equation}
Collision preserves zeroth and first moments; periodic streaming preserves the global sum of populations.
The relaxation parameter $\tau_{\mathrm{LB}}$ has units of time. In common lattice units $\Delta t=c_\ell=1$,
the relaxation rate is $s=1/\tau_{\mathrm{LB}}$. Boundary reflection or reconstruction must also account for mass and momentum exchange.

\section{Chapman--Enskog Expansion and Viscosity}\index{Chapman--Enskog expansion}
Write $D_i=\partial_t+\bm c_i\cdot\nabla$.
A smooth long-wave Taylor expansion of~\eqref{eq:lbe} gives
\[
D_if_i+\frac{\Delta t}{2}D_i^2f_i+\cdots
=-\frac{f_i-f_i^{\mathrm{eq}}}{\tau_{\mathrm{LB}}}.
\]
Order the spatial and temporal derivatives by a small parameter and write
$f_i=f_i^{(0)}+\epsilon f_i^{(1)}+\epsilon^2f_i^{(2)}+\cdots$,
Impose zero mass and momentum moments on the nonequilibrium parts.
Zeroth order gives $f_i^{(0)}=f_i^{\mathrm{eq}}$; first order gives
\[
f_i^{(1)}=-\tau_{\mathrm{LB}}D_i^{(1)}f_i^{\mathrm{eq}}.
\]
Zeroth and first moments recover the leading isothermal Euler system;
the second moment converts $f_i^{(1)}$ into viscous stress, while the Taylor term supplies a time-discretization correction.
On the low-Mach, nearly incompressible viscous scale,
\[
\partial_t\rho+\nabla\cdot(\rho\bm u)=0,
\]
\[
\rho_0(\partial_t\bm u+\bm u\cdot\nabla\bm u)
=-\nabla p+\rho_0\nu\Delta\bm u,\qquad
\nabla\cdot\bm u=0
\]
is the corresponding limiting equation. Here $p$ is the small pressure perturbation associated with $c_s^2\rho$,
$\rho/\rho_0=1+O(\mathrm{Ma}^2)$, and
\begin{equation}\label{eq:lbm-viscosity}
\nu=c_s^2\left(\tau_{\mathrm{LB}}-\frac{\Delta t}{2}\right).
\end{equation}
Positive viscosity requires $\tau_{\mathrm{LB}}>\Delta t/2$.
Positive viscosity sets a basic relaxation-parameter range; stability of the complete update also involves equilibrium positivity, boundaries, and nonhydrodynamic modes.
The relation between continuous BGK collision time and discrete $\tau_{\mathrm{LB}}$ depends on time integration;
equation~\eqref{eq:lbm-viscosity} applies to the update specified here.
The macroscopic limit of this standard second-order isothermal model assumes low Mach number and long wavelengths~\cite{qian1992}.

\aiexercise{List all nine D2Q9 velocities and weights, and compute equilibrium mass, momentum, and second moments directly. Write collision and streaming on a periodic lattice without physical boundaries, and prove global mass conservation step by step.}
{At fixed target kinematic viscosity, change $\Delta t$ and recompute $\tau_{\mathrm{LB}}$. For a periodic shear wave $u_x(y,0)=U_0\sin ky$, measure the decay rate $\nu k^2$ and distinguish lattice from physical units.}

\chapter{Modern Lattice Boltzmann Models and Algorithms}\label{ch:16}
\begin{chaptersummary}
Single-relaxation isothermal LBM established the route from discrete velocities to macroscopic transport. This chapter compares collision designs through moments, cumulants, and discrete entropy, then identifies additional variables and moment conditions for thermal, compressible, and boundary models. Model selection depends on target physics, Mach number, temperature variation, and a verified stability range.
\end{chaptersummary}
\learninggoals{Distinguish raw moments, central moments, and cumulants; write MRT and discrete entropy constraints; explain additional moment requirements for thermal and compressible models; design boundary and Galilean-error tests.}
\studyroute{Model structures and basic formulas are given; explicit MRT tables, high-order cumulant back-transforms, and coefficients of particular compressible lattices are omitted and should be obtained from the selected original model.}{MRT lattice Boltzmann; cumulant collision; entropic lattice Boltzmann; thermal lattice; Galilean invariance}{Compare the conserved moments and adjustable nonconserved modes of BGK, MRT, and cumulant collision, using a common viscosity and error measures for the same shear wave.}

\section{MRT: Relaxation in Moment Space}\index{MRT--LBM}
For an invertible moment transform $m=Tf$, multiple-relaxation-time collision is
\[
m^*=m-S(m-m^{\mathrm{eq}}),\qquad
f^*=T^{-1}m^*.
\]
Conserved moments have zero rates in $S=\diag(s_a)$; other modes may relax at different rates.
The shear-stress rate $s_\nu$ determines
\[
\nu=c_s^2\Delta t\left(\frac1{s_\nu}-\frac12\right).
\]
Other nonconserved rates affect acoustics, bulk viscosity, and lattice-scale damping.
For local diagonal relaxation, $0<s_a<2$ controls individual relaxation modes;
the full update also includes streaming, equilibrium linearization, and boundaries, so its combined spectrum must be checked~\cite{dhumieres2002}.

\conceptbridge{Why specify the moments in a multiple-relaxation model rather than list relaxation rates alone?}
{The moment transform identifies conserved, stress, and higher-order nonequilibrium modes.
Applying identical rates to different moments yields different collision operators.
Central-moment and cumulant collisions also change the dependence of the collision coordinates on local velocity.}
{Collision coordinates $\longrightarrow$ modal relaxation $\longrightarrow$ transport and nonhydrodynamic errors.}

\section{Central Moments and Cumulants}
Raw moments are $\sum_i f_i\bm c_i^\alpha$; central moments replace the velocity monomial with
$(\bm c_i-\bm u)^\alpha$.
For $\rho>0$, define a generating function using normalized populations $f_i/\rho$:
\[
K(\bm t)=\log\left(\frac1\rho\sum_i f_i
\exp[\bm t\cdot(\bm c_i-\bm u)]\right),\qquad
\kappa_\alpha=\left.\partial_{\bm t}^{\alpha}K\right|_{\bm t=0}.
\]
Second cumulants give the central covariance; higher cumulants subtract products of lower-order moments.
Cumulant collision relaxes these coordinates and reconstructs $f_i$,
controlling coupling among different orders of nonequilibrium information~\cite{geier2015}.
Finite velocity sets still constrain independently representable moments. Equilibrium, rates, and corrections must be specified together.

\section{Discrete Entropy and Positivity}\index{Entropic lattice Boltzmann}
For positive populations, define the convex function
\[
H(f)=\sum_i f_i\log\frac{f_i}{w_i}.
\]
The entropic equilibrium minimizes $H$ under mass and momentum constraints.
One entropic collision searches along $\Delta=f^{\mathrm{eq}}-f$ for a nonzero positivity-preserving root
$H(f+\alpha\Delta)=H(f)$, then uses $0<\beta\le1$ to set
$f^*=f+\alpha\beta\Delta$.
Convexity gives $H(f^*)\le H(f)$ along this segment.
This entropy estimate assumes an entropic equilibrium, a positivity domain, and a nonzero root. Polynomial equilibria or fixed BGK rates require the same conditions to be established for their collision maps.
For the original discrete entropy construction, see~\cite{karlin1999}.

\section{Thermal Models, Compressible Models, and Boundaries}
The isothermal relation $p=c_s^2\rho$ fixes the corresponding equation of state.
Heat transport requires an additional energy distribution or a richer velocity set,
with matched heat-flux moments, specific heats, and target Prandtl number.
Compressible flow additionally requires the correct energy flux, variable-temperature response, and sufficient velocity moments;
the macroscopic limit must be verified for the selected velocity set, equilibrium expansion, and collision model.

Halfway bounce-back at a stationary wall reverses outgoing populations into incoming ones;
wall location and slip error depend on geometric treatment. Moving walls require momentum-exchange corrections,
and inlet/outlet models reconstruct unknown incoming populations from prescribed macroscopic data.
Thermal boundaries additionally specify energy exchange. Curved-wall interpolation, diffuse reflection, and rarefied-gas boundaries have distinct validity conditions.

\section{Comparing Stability and Galilean Errors}
Superpose a uniform background velocity $U_b$ on a shear wave. Its continuum decay rate should be independent of $U_b$.
Viscosity and phase comparisons across background velocities measure Galilean errors due to higher lattice moments.
Record Mach number, viscosity, lattice resolution, wall location, and population positivity.
Linear spectra, discrete entropy, and nonlinear experiments probe different stability and error mechanisms.

\aiexercise{Choose an explicit invertible D2Q9 moment transform and identify conserved and shear modes. Set all nonconserved rates equal and determine when the associated equilibrium moments recover the BGK update.}
{For a periodic shear wave at two background velocities, compare decay errors of BGK, MRT, and a documented cumulant model. Identify the additional state variables and moment constraints needed for temperature variation.}

\part{Reduction and Closure in Complex Physical Systems}
\chapter{Rarefied Gas Dynamics and Nonequilibrium Closure}\label{ch:17}
\begin{chaptersummary}
Velocity moments connected kinetics to continuum models in Chapters 1 and 2. Stronger nonequilibrium calls for reconsidering the retained state and closure. Knudsen-number comparisons organize Chapman--Enskog, Grad, and simplified collision models, followed by DSMC's stochastic-particle approximation of Boltzmann dynamics. Model, closure, and sampling errors can then be tracked separately.
\end{chaptersummary}
\learninggoals{Distinguish Chapman--Enskog and finite-moment closure; identify the variables in Grad 13-moment and ES--BGK models; assess continuum validity using local scales; identify DSMC sampling error.}
\studyroute{Hyperbolicity proofs for higher-moment systems, general H theorems, and full DSMC collision-sampling constructions are omitted. Transport coefficients depend on the chosen collision kernel.}{Grad 13 moment; Chapman--Enskog; ES BGK; Knudsen layer; DSMC}{For planar Couette flow, compare Newtonian stress, Grad stress, and the kinetic distribution. Identify what boundary data and bulk closure each omit in a Knudsen layer.}

\section{Knudsen Number and Nonequilibrium Information}\index{Rarefied gases}
The ratio $\mathrm{Kn}=\lambda/L$ compares mean free path with a macroscopic length.
Conventional continuum, slip, transition, and free-molecular ranges are modeling guidelines;
local gradients, wall layers, and accuracy targets also determine applicability.
For example, $L_T=T/|\nabla T|$ gives a temperature-gradient Knudsen number $\lambda/L_T$.
Even for small global $\lambda/L$, shocks and wall regions may require nonequilibrium descriptions.

\section{Two Closure Routes: Asymptotic Manifolds and Moment Systems}
At small Knudsen number, the Chapman--Enskog expansion is
\[
f=M[n,\bm u,T]+\mathrm{Kn}\,f^{(1)}
+\mathrm{Kn}^2f^{(2)}+\cdots.
\]
Nonequilibrium corrections have zero conserved moments, so $n,\bm u,T$ are defined by Maxwellian parameters.
Velocity integration of the first correction gives Newtonian stress and Fourier heat flux;
their transport coefficients depend on the linearized collision operator.

Grad methods instead retain higher moments as independent evolving variables.
For a three-dimensional monatomic gas, the 13 moments are
$n$, $\bm u$, $T$, five traceless-stress components, and three heat-flux components.
A corresponding distribution representation has the structure
\[
f_{13}=M\left[1+
A_{ij}\left(c_ic_j-\tfrac13|\bm c|^2\delta_{ij}\right)
+B_i c_i(|\bm c|^2-5\theta)\right],
\qquad \theta=k_BT/m.
\]
Here $A_{ij}$ is symmetric and traceless, with $A,B$ determined by the prescribed stress and heat flux.
Substitution into higher moments gives a finite-moment closure~\cite{grad1949}.
Polynomial corrections can become negative in velocity tails, so realizability and hyperbolicity require explicit checks;
regularized moment models must be selected for their stated regimes.

\conceptbridge{Why do asymptotic and Grad closures differ if both describe stress and heat flux?}
{An asymptotic closure uses small Knudsen number to constrain the nonequilibrium distribution, expressing stress and heat flux through macroscopic gradients.
Grad models retain these moments as independent variables and close the higher moments in their evolution equations.
They therefore use different state spaces and time scales, with different errors and validity regions.}
{Small-parameter asymptotic relations / independent higher moments $\longrightarrow$ two closure families.}

\section{BGK and ES--BGK}
Continuous BGK relaxes toward a Maxwellian while preserving mass, momentum, and energy;
its single relaxation time links transport processes, giving $\mathrm{Pr}=1$ for the standard monatomic model.
ES--BGK uses an anisotropic Gaussian:
\[
G_\chi=\frac{n}{(2\pi)^{3/2}\sqrt{\det\mathcal T_\chi}}
\exp\left(-\frac12\bm c^{\mathsf T}\mathcal T_\chi^{-1}\bm c\right),
\]
\[
\mathcal T_\chi=(1-\chi)\theta I+\chi\Theta,\qquad
\Theta=\bm P/\rho,\qquad
C[f]=\frac{G_\chi-f}{\tau_c}.
\]
The identity $\tr\Theta=3\theta$ preserves energy; the covariance must remain positive definite.
This parameterization allows $\mathrm{Pr}=1/(1-\chi)$;
a common monatomic choice is $\chi=-1/2$, giving $\mathrm{Pr}=2/3$,
with viscosity matched to the model's collision-frequency definition~\cite{holway1966}.
Changes in discrete velocities, quadrature, or time stepping require renewed checks of collision invariants and entropy.

\section{DSMC and Error Separation}
DSMC represents a distribution with weighted simulation particles, alternating free streaming and random collisions in small cells.
Collision partners and scattering angles are sampled from the Boltzmann kernel;
wall or inlet distributions generate incoming particles~\cite{bird1994}.
Finite particle counts introduce statistical fluctuations; cell size and time step introduce numerical bias.
Resolving mean free paths, collision times, and local gradients guides refinement.

To compare continuum and kinetic models, first fix collision physics,
then reduce mesh, time-step, and sampling errors separately.
This separates closure error from numerical error.

\aiexercise{Multiply a standard normal density by $1+\alpha H_3(v)$, with $H_3(v)=v^3-3v$. Explain why any nonzero $\alpha$ can produce negative velocity tails. Calculate zeroth, first, and second moments, and distinguish low-order moment matching from positivity.}
{Design Couette-flow comparisons varying Knudsen number, velocity-space resolution, and DSMC sample count separately. Measure slip velocity and stress, distinguishing closure, boundary-model, and sampling errors.}

\chapter{MHD and Plasma Fluid Models}\label{ch:18}
\begin{chaptersummary}
Maxwell and species-kinetic descriptions contain charge separation and electron and ion time scales. MHD uses low-frequency, quasineutral, and scale orderings to form a single-fluid model, closed by pressure, heat flux, and Ohm relations. This chapter gives conservative ideal MHD, then uses resistivity and pressure anisotropy to show how physical closure changes energy exchange and numerical constraints.
\end{chaptersummary}
\learninggoals{State ideal-MHD unknowns and assumptions; check total energy and magnetic divergence; distinguish electrical resistivity from magnetic diffusivity; explain pressure-tensor closure and its scales of validity.}
\studyroute{The full asymptotic derivation of generalized Ohm's law from two-fluid equations, transport-coefficient tables, and anisotropic-instability dispersion calculations are omitted.}{ideal MHD; resistive MHD; generalized Ohm law; gyrotropic pressure; CGL closure}{Derive total-energy flux from ideal-MHD momentum, induction, and internal-energy equations. Check magnetic pressure and tension, identifying every use of the magnetic-divergence constraint.}

\section{The Low-Frequency Single-Fluid State}\index{MHD}
Use SI units and single-fluid variables $\rho,\bm u,p,\bm B$.
Quasineutrality makes bulk charge density relatively small. Low-frequency ordering neglects displacement current, giving
\[
\bm J=\frac1{\mu_0}\nabla\times\bm B,\qquad
\partial_t\bm B=-\nabla\times\bm E.
\]
Ideal Ohm's law is $\bm E+\bm u\times\bm B=0$.
Its validity also requires Hall, electron-pressure-gradient, and electron-inertia terms to be relatively small;
the ion inertial length relative to the target length, for example, must satisfy the chosen ordering.
Assume isotropic pressure with ideal-gas closure $p=(\gamma-1)\rho e$.

\section{The Conservative Ideal-MHD System}
Define
\[
\mathcal E=\rho e+\tfrac12\rho|\bm u|^2+\frac{|\bm B|^2}{2\mu_0},
\qquad p_t=p+\frac{|\bm B|^2}{2\mu_0}.
\]
Without external forces, viscosity, or heat conduction,
\begin{equation}\label{eq:ideal-mhd}
\begin{aligned}
\partial_t\rho+\nabla\cdot(\rho\bm u)&=0,\\
\partial_t(\rho\bm u)+\nabla\cdot
\left(\rho\bm u\bm u+p_tI-\frac{\bm B\bm B}{\mu_0}\right)&=0,\\
\partial_t\bm B-\nabla\times(\bm u\times\bm B)&=0,\\
\partial_t\mathcal E+\nabla\cdot
\left[(\mathcal E+p_t)\bm u-\frac{(\bm u\cdot\bm B)\bm B}{\mu_0}\right]&=0,
\end{aligned}
\end{equation}
with $\nabla\cdot\bm B=0$.
The magnetic-pressure and tension form uses this constraint; its continuous propagation follows from $\nabla\cdot\nabla\times=0$.
A discrete design should verify cell conservation, magnetic divergence, positive density, and positive internal energy.

\conceptbridge{Why retain magnetic divergence and closure conditions separately after writing conservative MHD?}
{Conservation equations describe flux exchange among state variables; magnetic divergence restricts admissible magnetic fields.
Pressure and Ohm relations encode scale and transport assumptions. These conditions constrain states, evolution, and unresolved physics in different ways,
and therefore require distinct discrete verification measures.}
{Low-frequency reduction $\longrightarrow$ single-fluid conservation + magnetic constraint + physical closure.}

\section{Resistivity, Diffusion, and Energy Conversion}
For electrical resistivity $\eta_e\ge0$,
\[
\bm E+\bm u\times\bm B=\eta_e\bm J.
\]
if $\eta_e,\mu_0$ are constant and $\nabla\cdot\bm B=0$,
\[
\partial_t\bm B=\nabla\times(\bm u\times\bm B)
+\eta_m\Delta\bm B,\qquad \eta_m=\eta_e/\mu_0.
\]
$\eta_e$ has units $\Omega\cdot\mathrm m$, while $\eta_m$ has units $\mathrm m^2/\mathrm s$.
Magnetic energy loses $\eta_e|\bm J|^2$ through Ohmic dissipation, and internal energy gains the same Joule heat;
total energy also requires a Poynting flux compatible with the resistive electric field.
Explicit magnetic diffusion typically imposes an $O(h^2/\eta_m)$ step restriction. Propagation and diffusion scales can motivate IMEX treatment.

\section{Pressure Anisotropy and Transport Closure}
In a strong magnetic field, write
\[
\bm P=p_\perp I+(p_\parallel-p_\perp)\bm b\bm b+\bm\Pi_{\mathrm{ng}},
\qquad \bm b=\bm B/|\bm B|.
\]
Here $\bm\Pi_{\mathrm{ng}}$ is nongyrotropic; heat flux may also be decomposed into parallel and perpendicular components.
In the collisionless double-adiabatic approximation, neglecting heat flux and finite-Larmor-radius effects,
\[
\frac{D}{Dt}\left(\frac{p_\perp}{\rho|\bm B|}\right)=0,\qquad
\frac{D}{Dt}\left(\frac{p_\parallel|\bm B|^2}{\rho^3}\right)=0.
\]
The CGL relations connect parallel and perpendicular pressures to density and magnetic field, closing pressure under the double-adiabatic ordering.
Collisional transport and nonlocal heat flux require corresponding models and may change well-posedness and instability thresholds~\cite{cgl1956}.

\aiexercise{For constant-resistivity MHD in a periodic domain, calculate magnetic-energy loss and internal-energy Joule heating separately. Prove cancellation of the volume conversion terms and write the boundary energy flux.}
{Superpose a small Alfv\'en wave on an equilibrium magnetic field and measure wave speed, phase, and discrete magnetic divergence. Introduce parallel and perpendicular pressures and identify the equations that must change from isotropic closure.}

\chapter{Gyrokinetics and Multiscale Reduction}\label{ch:19}
\begin{chaptersummary}
The strong-field ordering of Chapter 4 separates fast gyromotion from slow transport. Coordinate transformations and ordered averaging remove the fast angle while retaining parallel-velocity and magnetic-moment distributions. Gyroaveraging and reconstruction of physical moments reconnect reduced particles with self-consistent fields. A uniform-field electrostatic model makes this five-dimensional reduction concrete.
\end{chaptersummary}
\learninggoals{Specify strong-field orderings and five-dimensional coordinates; distinguish guiding centers from gyrocenters; calculate Fourier gyroaverages; separate kinetic reduction, moment closure, field coupling, and numerical discretization.}
\studyroute{A controlled uniform-field electrostatic example is given; general electromagnetic Lie transforms, second-order Hamiltonians, magnetization, and collision transformations in magnetic geometry are omitted. A complete model must use a consistent variational order.}{guiding center; gyrocenter; gyroaverage; finite Larmor radius; gyrokinetic polarization; electromagnetic gyrokinetics}{Gyroaverage a perpendicular Fourier potential in a uniform field, derive the Bessel factor, and write linear quasineutrality for a specified nonadiabatic distribution variable.}

\section{Scale Ordering and the Five-Dimensional State}\index{Gyrokinetics}
Using equilibrium scale $L_B$ and ion thermal gyroradius $\rho_{Li}$, a common perturbation ordering is
\[
\epsilon=\rho_{Li}/L_B\ll1,\qquad
\omega/\Omega_i=O(\epsilon),\qquad
k_\parallel/k_\perp=O(\epsilon),\qquad
k_\perp\rho_{Li}=O(1).
\]
Weak perturbations may additionally satisfy $|q_i\delta\phi|/T_i=O(\epsilon)$.
Below, $T_s$ denotes temperature in energy units, corresponding to $k_BT_s$ earlier.
Electromagnetic, collisional, and electron-scale models require additional orderings.

The original phase space $(\bm x,\bm v)$ can be transformed to
$(\bm R,v_\parallel,\mu,\vartheta)$,
where $\mu=m v_\perp^2/(2B_0)$ and $\vartheta$ is the fast gyrophase.
A guiding-center transformation first handles rapid gyromotion in the equilibrium field;
a gyrocenter transformation then handles perturbation-induced angle dependence.
The reduced distribution $F_s(\bm R,v_\parallel,\mu,t)$ no longer depends independently on $\vartheta$,
retaining three position and two velocity coordinates.
For uniform $B_0$ and a gyrophase-independent function $A$, velocity moments use
\[
\int A\dd^3v=\frac{2\pi B_0}{m_s}
\int_{-\infty}^{\infty}\int_0^\infty A\dd\mu\dd v_\parallel.
\]
The angle integral has been performed analytically; parallel velocity and magnetic moment remain independent variables.

\conceptbridge{Why do field equations still require polarization and gyroaveraging after a fast angle has been removed?}
{The position coordinate changes from particle position to gyrocenter, but density and current are measured at physical positions.
Pushing the reduced distribution back to physical space produces finite-Larmor-radius, polarization, and magnetization contributions.
These arise at the same transformation order as the reduced dynamics and determine consistent fields and energy pairings.}
{Phase-space transformation $\longrightarrow$ five-dimensional evolution $\longrightarrow$ physical moments and self-consistent fields.}

\section{Gyroaveraging and Finite Larmor Radius}
Write particle position as $\bm x=\bm R+\bm\rho(\vartheta)$.
At fixed gyrocenter and magnetic moment, define
\[
\langle\phi\rangle_{\bm R}
=\frac1{2\pi}\int_0^{2\pi}
\phi(\bm R+\bm\rho(\vartheta))\dd\vartheta.
\]
For a plane wave in a uniform magnetic field,
\[
\left\langle\ee^{\ii\bm k\cdot\bm x}\right\rangle_{\bm R}
=\ee^{\ii\bm k\cdot\bm R}
J_0(k_\perp v_\perp/\Omega_s).
\]
finite $k_\perp\rho_{Ls}$ changes the particle response through ring averaging.
The small parameter removes the fast angle, while the Bessel factor retains its influence on slow-scale coupling.

\section{Kinetic Structure and a Specific Model}
General gyrocenter dynamics can be written in characteristic form:
\[
\partial_tF_s+\dot{\bm R}_s\cdot\nabla_{\bm R}F_s
+\dot v_{\parallel s}\partial_{v_\parallel}F_s=C_s^{\mathrm{GK}}[F].
\]
Characteristics, phase-space measure, and collisions must use the same reduced coordinates;
a nonuniform field requires the corresponding phase-space Jacobian in conservative form.

As a computable example, take uniform $B_0\bm b$, no equilibrium gradients, no collisions, and low-frequency electrostatic perturbations.
Let $F_{0s}$ be a uniform Maxwellian at temperature $T_s$,
and $h_s$ the usual nonadiabatic distribution, with its convention fixed by the physical-density formula below.
A corresponding nonlinear electrostatic system is
\begin{equation}\label{eq:slab-gk}
\partial_th_s+v_\parallel\partial_\parallel h_s
+\bm v_{E,s}\cdot\nabla_\perp h_s
=\frac{q_sF_{0s}}{T_s}\partial_t\langle\phi\rangle_{\bm R},
\qquad
\bm v_{E,s}=\frac{\bm b\times\nabla\langle\phi\rangle_{\bm R}}{B_0}.
\end{equation}
The drift is the $\bm E\times\bm B/B_0^2$ velocity of the ring-averaged electric field.
A uniform equilibrium has no gradient drive; density or temperature gradients require matching drive terms.
For general electromagnetic construction and energy conservation, see~\cite{brizard2007,abel2013}.

\section{Quasineutrality, Polarization, and Amp\`ere Coupling}
For a Fourier mode of this uniform-field model, the physical density perturbation is
\[
\delta n_{s,\bm k}
=-\frac{q_sn_{0s}}{T_s}\phi_{\bm k}
+\int J_{0s}h_{s,\bm k}\dd^3v,\qquad
J_{0s}=J_0(k_\perp v_\perp/\Omega_s).
\]
Thus quasineutrality, $\sum_sq_s\delta n_s=0$, gives
\begin{equation}\label{eq:gk-qn}
\sum_s\frac{q_s^2n_{0s}}{T_s}\phi_{\bm k}
=\sum_sq_s\int J_{0s}h_{s,\bm k}\dd^3v.
\end{equation}
If instead $g_{s,\bm k}=h_{s,\bm k}-(q_sF_{0s}/T_s)J_{0s}\phi_{\bm k}$,
the same relation becomes
\[
\sum_s\frac{q_s^2n_{0s}}{T_s}(1-\Gamma_{0s})\phi_{\bm k}
=\sum_sq_s\int J_{0s}g_{s,\bm k}\dd^3v,
\]
\[
\Gamma_{0s}=I_0(b_s)\ee^{-b_s},\qquad
b_s=k_\perp^2T_s/(m_s\Omega_s^2).
\]
The two forms use different distribution variables, so the polarization term appears in different places.
The zero-wave-number potential gauge and particle-number compatibility require separate specification.

Electromagnetic models also determine $A_\parallel$ through parallel current
and include $\delta B_\parallel$, polarization, and magnetization according to their ordering.
The Amp\`ere current must use a push-back consistent with the gyrocenter coordinates, distribution variable, and field equations.
These relations couple the retained velocity distributions to electromagnetic potentials; truncating higher moments is a separate closure choice.

\section{Collisions, Turbulent Transport, and Discretization Levels}
Gyrokinetic collision models must preserve the appropriate particle, momentum, and energy moments
and be compatible with the reduced measure and entropy dissipation. Turbulent particle and heat fluxes arise from correlations between distribution perturbations and radial drifts;
a particle flux has the structure $\langle\int\delta f_s\,v_{E,r}\dd^3v\rangle$.
Specific fluxes depend on the push-back definition and averaging region.

The five-dimensional state still requires discretization in position, parallel velocity, magnetic moment, and time.
Domain extent, velocity-tail truncation, and collision resolution control different errors;
spectral, DG, or particle representations are subsequent numerical choices.

\aiexercise{For a Maxwellian equilibrium and uniform field, calculate $\langle\exp(\ii\bm k_\perp\cdot\bm\rho)\rangle$ and verify $\int F_{0s}J_{0s}^2\dd^3v=n_{0s}\Gamma_{0s}$. Prove that the $h_s$ and $g_s$ quasineutrality forms are equivalent.}
{Starting from~\eqref{eq:slab-gk}, fix periodic boundaries and velocity truncation, and list the data needed to check quasineutrality residuals, free energy, and velocity-space convergence. For further reading, choose one electromagnetic model and map its distribution variable, Jacobian, collision term, and field equations consistently.}

\part{Meshes and Verification in Computational Physics}
\chapter{Meshes as the Interface between Geometry and Discrete Operators}\label{ch:20}
\begin{chaptersummary}
After model and space selection, the mesh introduces physical geometry, boundaries, and local scales into the operator. Reference-element mappings explain how shape and size change gradients, integration, and interface pairings, leading to curved geometry and conservation checks. These compatibility requirements provide the common foundation for the next chapter's mesh techniques.
\end{chaptersummary}
\learninggoals{Transform gradients and integrals from reference to physical elements; distinguish geometric and approximation errors; explain how shape, adjacency, and boundary representation affect spatial operators.}
\studyroute{Higher-dimensional mesh generation, geometric-approximation error proofs, and matrix implementation are omitted. This chapter retains operator structure; algebraic solvers and parallel data layout are developed in Volume~0.}{element mapping; mesh quality; curved element; metric identity; discrete geometric conservation}{Compare Jacobians and linear finite-element stiffness matrices for an equilateral and a slender triangle, identifying the geometric factor that amplifies gradients.}
\section{Geometry, Adjacency, and Discrete Coupling}\index{Mesh geometry}
Meshes determine geometric representation error, local discretization error, stability restrictions, and matrix nonzero patterns. Geometry handling and mesh-quality control are central computational tasks.

\begin{authornote}
\textbf{Potential Question:}
"Why does a mesh affect matrix structure, stability restrictions, and solver behavior as well as geometric shape?"

\textbf{Author Supplement:}
A mesh defines not only sampling locations but also information exchange between elements. Adjacency determines discrete coupling; element shapes and sizes enter derivatives, integrals, and time-step estimates; boundary representation affects boundary-condition enforcement. The same continuous equation therefore produces different couplings and errors on different meshes.

\textbf{Core Conceptual Connection:}
\[
\text{mesh geometry}
\longrightarrow
\text{discrete coupling}
\longrightarrow
\text{matrix and solver behavior}
\]
\end{authornote}

A body-fitted mesh aligns its boundary with the physical surface. Boundary conditions are imposed directly, and wall layers can be resolved, but complex geometry and large motion complicate generation and quality control.

\section{Reference Mappings and Geometric Scales}
Let $x=F_K(\xi)$ with invertible $A_K=\partial F_K/\partial\xi$. Then
\[
\nabla_xu=A_K^{-\mathsf T}\nabla_\xi\widehat u,\qquad
\dd x=|\det A_K|\dd\xi.
\]
A Poisson stiffness entry becomes
\[
K_{ij}^{(K)}
=\int_{\widehat K}(\nabla_\xi\phi_i)^{\mathsf T}
A_K^{-1}A_K^{-\mathsf T}\nabla_\xi\phi_j\,|\det A_K|\dd\xi.
\]
Slender or degenerate elements change mapping singular values, amplifying directional gradients and error constants.
Structured meshes have regular connectivity; unstructured meshes permit more flexible shapes and adjacency.
Both require boundary orientations, face measures, and discrete integration rules.

\section{Curved Elements and Geometric Error}
High-order field approximation often requires correspondingly accurate curved geometry.
Mapping degree and volume and surface quadrature jointly affect total error.
For conservation laws, preservation of constants requires the discrete geometric factors to satisfy the relevant metric identities;
shared face geometry and opposite normals are the basis of flux cancellation.
Begin verification with constant fields, low-order polynomials, and analytic surfaces to check geometric integrals, then assess convergence of physical solutions.

\aiexercise{For $x=A\xi+b$, choose rotational, isotropic-scaling, and strongly stretched matrices $A$. Derive gradient, mass, and stiffness scalings, and compare errors for the same linear field on the three meshes.}
{Represent a circular boundary with straight segments and quadratic curves. At fixed field degree, design convergence experiments separating geometric error from spatial approximation error.}

\chapter{Complex Geometry and Mesh Techniques}\label{ch:21}
\begin{chaptersummary}
Fixed body-fitted meshes represent object boundaries directly; complex motion and changing local scales require more flexible geometry. Immersed boundaries, tree meshes, ALE, and overset grids connect geometry to evolution through forcing, common subfaces, moving control volumes, and interpolation. Conservative transfer and geometric conservation provide common principles across these layouts.
\end{chaptersummary}
\learninggoals{Explain the states and interfaces of four mesh techniques; derive relative fluxes for moving control volumes; state conservation requirements for coarse--fine transfers, common subfaces, and overset interpolation.}
\studyroute{Geometry, interface transfer, and conservation organize the mesh equations. Full mesh generators, hole-cutting searches, and boundary-forcing implementations are left for further study.}{immersed boundary; quadtree octree; AMR refluxing; ALE; geometric conservation law; overset}{List shared fluxes for one coarse face adjoining two fine faces, then test a constant solution in a moving control volume and compare spatial with geometric conservation.}
\section{Immersed Boundaries and Background Meshes}
An immersed-boundary method places complex objects on a simple background mesh without requiring boundary alignment. Boundary action enters the momentum equation:
\[
\rho\left(
\frac{\partial\bm u}{\partial t}
+\bm u\cdot\nabla\bm u
\right)
=
-\nabla p+\mu\nabla^2\bm u+\bm f_{\mathrm{IB}}.
\]

\begin{authornote}
\textbf{Potential Question:}
"Without a boundary-fitted mesh, why can \(\bm f_{\mathrm{IB}}\) represent the object's effect?"

\textbf{Author Supplement:}
A body-fitted mesh represents the object surface directly as a mesh boundary where conditions can be imposed. An immersed surface cuts through the background mesh, so added forces or modified nearby discrete equations must enforce its effect. This permits more flexible geometry handling, while requiring explicit checks of near-wall accuracy and conservation.

\textbf{Core Conceptual Connection:}
\[
\text{body-fitted boundary condition}
\longrightarrow
\text{immersed forcing}
\longrightarrow
\text{background-grid computation}
\]
\end{authornote}

Continuous forcing transfers forces between Lagrangian boundary points and an Eulerian grid through smooth kernels; discrete forcing modifies nearby discrete equations directly.

Regular background grids simplify computation and complex motion, while near-wall accuracy, flux conservation, and thin structures require corresponding boundary discretizations and checks.

\section{Quadtrees, Octrees, and AMR}
Quadtrees recursively divide a two-dimensional cell into four children; octrees divide a three-dimensional cell into eight. Resolution is concentrated near shocks, shear layers, vortices, interfaces, or high-error regions.

\begin{authornote}
\textbf{Potential Question:}
"Why does local refinement introduce 2:1 balancing, nonconforming interfaces, and conservative transfers?"

\textbf{Author Supplement:}
Refinement changes more than resolution: it creates new coarse--fine adjacency, requiring compatible fluxes across differently sized faces. Moving refinement regions also require new spatial projections and interface fluxes so that transfers between old and new degrees of freedom remain conservative.

\textbf{Core Conceptual Connection:}
\[
\text{local refinement}
\longrightarrow
\text{nonmatching interfaces}
\longrightarrow
\text{conservative transfer}
\]
\end{authornote}

Common implementation constraints include
\begin{itemize}
  \item maintaining 2:1 level balance between neighbors;
  \item treating coarse--fine nonconforming interfaces;
  \item conservatively transferring variables during refinement and coarsening;
  \item rebuilding interface projections and flux matching after adaptation.
\end{itemize}

\section{Moving Meshes, ALE, and Geometric Conservation}
In an arbitrary Lagrangian--Eulerian description, the mesh moves with velocity \(\bm w\). Conservation becomes
\[
\frac{\dd}{\dd t}\int_{\Omega(t)}U\dd V
+
\int_{\partial\Omega(t)}
[\bm F(U)-U\bm w]\cdot\bm n\dd S
=0.
\]

\begin{authornote}
\textbf{Potential Question:}
"Why does mesh motion change the flux from \(\bm F(U)\) to \(\bm F(U)-U\bm w\)?"

\textbf{Author Supplement:}
A moving boundary changes the amount inside a control volume both through physical transport and by sweeping through space. Integral conservation therefore involves a relative flux combining physical flux and mesh velocity. Discrete geometric conservation requires volume changes to agree with swept volumes; otherwise even a uniform solution may change artificially.

\textbf{Core Conceptual Connection:}
\[
\text{fixed control volume}
\longrightarrow
\text{moving control volume}
\longrightarrow
\text{ALE flux}
\]
\end{authornote}

A moving mesh must satisfy geometric conservation: mesh motion alone must not alter a constant physical solution.

Mesh motion can use spring analogies, Laplace smoothing, linear elasticity, radial basis functions, or local reconstruction. Large deformations may require remeshing.

\section{Overset Grids and Nonconforming Coupling}
Overset grids superpose independent blocks, such as background, body-fitted, and moving local grids. Donor cells supply interpolated data to receiver cells.

\begin{authornote}
\textbf{Potential Question:}
"How do overlapping grids define one computational domain? Which structures can interpolation compromise?"

\textbf{Author Supplement:}
Several easily generated or moved subgrids jointly cover the domain. Covered cells must be identified to avoid duplicate evolution, while donor-to-receiver interpolation couples grid boundaries. Ordinary interpolation need not balance the fluxes on both sides, so global conservation requires additional design or verification.

\textbf{Core Conceptual Connection:}
\[
\text{single conforming mesh}
\longrightarrow
\text{overlapping mesh blocks}
\longrightarrow
\text{interpolation coupling}
\]
\end{authornote}

Independent component meshes facilitate large relative rigid-body motion. Challenges include hole cutting, donor search, conservative interpolation, and dynamic connectivity.


\section{Consistency in Space, Time, and Geometry}\index{Geometric conservation law}
Mortars can organize common subfaces at coarse--fine interfaces. Overset grids additionally require a well-defined active physical coverage
to avoid double-counting conserved quantities. State interpolation, flux matching, and conservative remapping perform different tasks.

For a moving element, the continuous geometric identity is
\[
\frac{\dd |K(t)|}{\dd t}
=\int_{\partial K(t)}\bm w\cdot\bm n\dd S.
\]
Discrete volume changes and mesh-swept volumes in ALE fluxes must use compatible time integration.
Substituting $U=U_0$ into the moving-control-volume update tests preservation of a uniform field.

\aiexercise{Let the endpoints of a one-dimensional element move at prescribed velocities. Derive the constant-solution update from the ALE integral form and state the identity required between volume integration and mesh flux.}
{Draw a tree-mesh coarse--fine interface and two overlapping grids. Label state transfers, active control volumes, and shared fluxes, then explain how global mass change is measured in each layout.}

\chapter{Method Selection and the Complete Computational Workflow}\label{ch:22}
\begin{chaptersummary}
Representation, scale reduction, closure, discretization, and meshes now form a complete modeling chain. Quantities of interest organize these choices, while separate tests address model, reduction, spatial, temporal, and algebraic errors. Conservation, stability, and convergence become an executable workflow, with method combinations determined by target accuracy, structural preservation, and verification evidence.
\end{chaptersummary}
\learninggoals{Select models and discretizations for a physical objective; distinguish five error classes; vary spatial resolution, time steps, and tolerances independently; distinguish numerical verification from physical validation.}
\studyroute{Verification strategies are organized here; experimental uncertainty, error estimators, and project-scale data management are omitted. Algebraic solvers and computer architecture are developed in Volume~0.}{verification and validation; grid convergence; temporal convergence; model error; reduction error; quantity of interest}{For a wave or transport problem, define a quantity of interest and vary model, mesh, degree, time step, and algebraic tolerance independently to separate error sources.}
\section{Method Selection and Combined Strategies}
For fixed geometry and high wall-accuracy requirements, body-fitted meshes directly resolve boundary layers.

For complex moving geometry, immersed boundaries simplify meshing when wall shear is secondary to the main target.

For strongly varying local scales, use quadtree or octree adaptation.

For moderate boundary motion, use moving meshes.

For large relative motion of several rigid bodies, use overset grids.

Industrial software often combines
\[
\text{Body-fitted wall layers}
+
\text{Adaptive far field},
\]
\[
\text{Octree background mesh}
+
\text{Immersed boundaries},
\]
\[
\text{Overset grids}
+
\text{Local moving meshes}.
\]

\begin{askai}
Ask AI:
\begin{quote}
For rotating blades with far-field sound propagation, compare the subproblems addressed by body-fitted, moving, overset, and adaptive far-field meshes. Propose a combined strategy based on each region's requirements.
\end{quote}

Follow-up:
\begin{quote}
Construct a selection table using geometric stationarity, motion amplitude, boundary-layer accuracy, local scale variation, and stable time step.
\end{quote}
\end{askai}
\section{The Complete Computational Workflow}
A reliable calculation generally proceeds through the following steps:

\begin{enumerate}
  \item Select a model using ratios such as \(\mathrm{Kn}\), \(\mathrm{Re}\), and \(\mathrm{Ma}\).
  \item Specify conserved variables, material laws, and boundary conditions.
  \item Choose a mesh strategy adapted to geometry and local scales.
  \item Choose discretizations meeting conservation, stability, and wave-accuracy requirements.
  \item Form the algebraic system through explicit or implicit time integration.
  \item Record algebraic tolerances and target-quantity sensitivity to residuals; solution algorithms are developed in Volume~0.
  \item Monitor model, mesh, and discretization errors, and linear and nonlinear residuals separately.
\end{enumerate}

Residuals measure satisfaction of algebraic equations. Model validity and discretization accuracy additionally require benchmarks, mesh and temporal convergence studies, and comparison with physical data.

\section{Five Error Classes and Verification Evidence}\index{Verification and validation}
\begin{center}\small
\begin{tabularx}{\textwidth}{@{}l>{\raggedright\arraybackslash}X@{}}
\toprule
Error level & Executable check\\
\midrule
Physical model & Compare with experiments, measurements, or a higher-fidelity model in its valid regime; report parameter and data uncertainties.\\
Reduction and closure & Vary the small parameter and compare full and reduced models; check unresolved moments and scale-dependent terms.\\
Spatial discretization & With sufficiently small time steps and tolerances, refine $h$ and increase $p$ separately; compare target quantities and expected orders.\\
Time discretization & With adequately resolved space, halve $\Delta t$ and examine order and accumulated phase error.\\
Algebraic solution & Tighten residual tolerances until target-quantity changes fall below spatial and temporal errors.\\
\bottomrule
\end{tabularx}
\end{center}
Geometric approximation, velocity-domain truncation, stochastic sampling, and boundary models should also be recorded at their respective levels.
Relating algebraic residuals to solution error requires operator stability or target-quantity sensitivity.

\section{Conservation, Stability, and Convergence Experiments}
For periodic or closed source-free systems, record total mass, momentum, and energy;
with sources and physical boundaries, record balance defects including inflow, outflow, and dissipation.
Energy estimates or discrete entropy specify admissible evolution; wave problems additionally require phase and group-velocity measurements.
Electromagnetic and MHD calculations should also record discrete divergence.

Smooth analytic solutions test expected order; manufactured solutions provide sources and boundaries exercising all operators.
Discontinuities, wall layers, and multiscale problems require error measures appropriate to their regularity and target quantities.
Numerical verification measures solution error for a given model; physical validation assesses how well that model represents the target phenomenon.

\conceptbridge{After residual, conservation, and grid-convergence checks pass, what evidence does the model still need?}
{Those checks assess mathematical behavior within a chosen model and discretization. Physical validation also requires measurements or appropriate benchmarks under the target conditions,
with uncertainty in constitutive laws, boundaries, and parameters. Together, these forms of evidence determine the credible scope of a computational conclusion.}
{Numerical verification $\longrightarrow$ model--data comparison $\longrightarrow$ credibility of the quantity of interest.}

\aiexercise{Build an error table for a periodic Alfv\'en wave or low-Mach shear wave, listing mesh, degree, time step, tolerance, conservation defect, and target-quantity error. Include controls that detect temporal error masking spatial convergence.}
{Change the target to wall heat flux or turbulent transport. Identify additional collision, constitutive, boundary, and statistical assumptions, and data that can test them independently.}

\section{The Dependency Chain}
\[
\begin{gathered}
\text{Liouville and statistical representation}
\longrightarrow\\
\text{Kinetic and continuum models}
\longrightarrow\text{Scale reduction and closure}
\end{gathered}
\]
\[
\begin{gathered}
\longrightarrow\text{Function spaces and weak forms}
\longrightarrow\\
\text{Residuals, fluxes, and high-order discretization}\\
\longrightarrow\text{Meshes and verification}.
\end{gathered}
\]
The retained state determines required closure; the continuous operator and target structures determine suitable residual constraints.
Target quantities and verification evidence then determine resolution and combinations of methods.

\backmatter
\cleardoublepage
\phantomsection
\addcontentsline{toc}{chapter}{Bibliography}
\begin{thebibliography}{99}
\raggedright
\bibitem{evans}
L. C. Evans, \emph{Partial Differential Equations}, 2nd ed., American Mathematical Society, 2010.
Further reading on weak derivatives, Sobolev spaces, energy estimates, and elliptic problems.
\bibitem{leveque}
R. J. LeVeque, \emph{Finite Volume Methods for Hyperbolic Problems}, Cambridge University Press, 2002.
Conservation laws, Riemann problems, and high-resolution fluxes.
\bibitem{hesthaven}
J. S. Hesthaven and T. Warburton, \emph{Nodal Discontinuous Galerkin Methods}, Springer, 2008.
Local bases, nodal representations, and energy and spectral analysis.
\bibitem{jiang1996}
G.-S. Jiang and C.-W. Shu,
\href{https://doi.org/10.1006/jcph.1996.0130}{\emph{Efficient Implementation of Weighted ENO Schemes}},
Journal of Computational Physics, 126 (1996), 202--228.
\bibitem{arnold2002}
D. N. Arnold, F. Brezzi, B. Cockburn, and L. D. Marini,
\href{https://doi.org/10.1137/S0036142901384162}{\emph{Unified Analysis of Discontinuous Galerkin Methods for Elliptic Problems}},
SIAM Journal on Numerical Analysis, 39 (2002), 1749--1779.
\bibitem{huynh2007}
H. T. Huynh,
\href{https://doi.org/10.2514/6.2007-4079}{\emph{A Flux Reconstruction Approach to High-Order Schemes Including Discontinuous Galerkin Methods}},
AIAA Paper 2007-4079, 2007.
\bibitem{vincent2011}
P. E. Vincent, P. Castonguay, and A. Jameson,
\href{https://doi.org/10.1007/s10915-010-9420-z}{\emph{A New Class of High-Order Energy Stable Flux Reconstruction Schemes}},
Journal of Scientific Computing, 47 (2011), 50--72.
\bibitem{ranocha2016}
H. Ranocha, P. \"Offner, and T. Sonar,
\href{https://arxiv.org/abs/1511.02052}{\emph{Summation-by-parts operators for correction procedure via reconstruction}},
2016, arXiv:1511.02052.
\bibitem{kopriva1996}
D. A. Kopriva,
\href{https://doi.org/10.1006/jcph.1996.0225}{\emph{A Conservative Staggered-Grid Chebyshev Multidomain Method for Compressible Flows. II. A Semi-Structured Method}},
Journal of Computational Physics, 128 (1996), 475--488.
\bibitem{dumbser2014}
M. Dumbser,
\href{https://arxiv.org/abs/1402.6897}{\emph{Arbitrary-Lagrangian-Eulerian ADER-WENO Finite Volume Schemes with Time-Accurate Local Time Stepping for Hyperbolic Conservation Laws}},
2014, doi:10.1016/j.cma.2014.07.019.
\bibitem{owren1992}
B. Owren and M. Zennaro,
\href{https://doi.org/10.1137/0913084}{\emph{Derivation of Efficient, Continuous, Explicit Runge--Kutta Methods}}, 1992.
\bibitem{altmann2010}
C. Altmann, G. Gassner, and C.-D. Munz,
\href{https://congress2.cimne.com/eccomas/proceedings/cfd2010/papers/01008.pdf}{\emph{An Explicit Space-Time Adaptive Discontinuous Galerkin Scheme}},
ECCOMAS CFD, 2010.
\bibitem{ascher1997}
U. M. Ascher, S. J. Ruuth, and R. J. Spiteri,
\href{https://doi.org/10.1016/S0168-9274(97)00056-1}{\emph{Implicit-explicit Runge--Kutta methods for time-dependent partial differential equations}},
Applied Numerical Mathematics, 25 (1997), 151--167.
\bibitem{qian1992}
Y. H. Qian, D. d'Humi\`eres, and P. Lallemand,
\href{https://doi.org/10.1209/0295-5075/17/6/001}{\emph{Lattice BGK Models for Navier-Stokes Equation}},
Europhysics Letters, 17 (1992), 479--484.
\bibitem{dhumieres2002}
D. d'Humi\`eres, I. Ginzburg, M. Krafczyk, P. Lallemand, and L.-S. Luo,
\href{https://doi.org/10.1098/rsta.2001.0955}{\emph{Multiple-relaxation-time lattice Boltzmann models in three dimensions}},
Philosophical Transactions of the Royal Society A, 360 (2002), 437--451.
\bibitem{geier2015}
M. Geier, M. Sch\"onherr, A. Pasquali, and M. Krafczyk,
\href{https://doi.org/10.1016/j.camwa.2015.05.001}{\emph{The cumulant lattice Boltzmann equation in three dimensions: Theory and validation}},
Computers \& Mathematics with Applications, 70 (2015), 507--547.
\bibitem{karlin1999}
I. V. Karlin, A. Ferrante, and H. C. \"Ottinger,
\href{https://doi.org/10.1209/epl/i1999-00370-1}{\emph{Perfect entropy functions of the Lattice Boltzmann method}},
Europhysics Letters, 47 (1999), 182--188.
\bibitem{grad1949}
H. Grad,
\href{https://doi.org/10.1002/cpa.3160020403}{\emph{On the kinetic theory of rarefied gases}},
Communications on Pure and Applied Mathematics, 2 (1949), 331--407.
\bibitem{holway1966}
L. H. Holway,
\href{https://doi.org/10.1063/1.1761920}{\emph{New Statistical Models for Kinetic Theory: Methods of Construction}},
Physics of Fluids, 9 (1966), 1658--1673.
\bibitem{bird1994}
G. A. Bird,
\href{https://doi.org/10.1093/oso/9780198561958.001.0001}{\emph{Molecular Gas Dynamics and the Direct Simulation of Gas Flows}},
Clarendon Press, 1994.
\bibitem{cgl1956}
G. F. Chew, M. L. Goldberger, and F. E. Low,
\href{https://static.ias.edu/pitp/2016/sites/pitp/files/chew_goldberger_low_1956-1.pdf}{\emph{The Boltzmann equation and the one-fluid hydromagnetic equations in the absence of particle collisions}},
Proceedings of the Royal Society A, 236 (1956), 112--118.
\bibitem{brizard2007}
A. J. Brizard and T. S. Hahm,
\href{https://doi.org/10.1103/RevModPhys.79.421}{\emph{Foundations of nonlinear gyrokinetic theory}},
Reviews of Modern Physics, 79 (2007), 421--468.
\href{https://bp-pub.pppl.gov/pub_report/2006/PPPL-4153.pdf}{PPPL preprint}.
\bibitem{abel2013}
I. G. Abel et al.,
\href{https://arxiv.org/abs/1209.4782}{\emph{Multiscale Gyrokinetics for Rotating Tokamak Plasmas: Fluctuations, Transport and Energy Flows}},
Reports on Progress in Physics, 76 (2013), 116201.
\end{thebibliography}

\clearpage
\phantomsection
\addcontentsline{toc}{chapter}{Terminology Index}
\printindex
\end{document}
